\makeatletter
\def\jobname{main_aux}
\makeatother
\documentclass[a4paper, 12pt]{article}

\usepackage[english]{babel}
\usepackage[utf8]{inputenc}
\usepackage{amsmath}
\usepackage{titling}
\usepackage{amssymb}
\usepackage{graphicx}
\usepackage[colorinlistoftodos]{todonotes}
\usepackage{setspace}
\usepackage[bf]{caption}
\usepackage[justification=centering]{caption}
\usepackage{pgfplots}
\pgfplotsset{compat=1.10}
\usepackage{tikz}
\usepackage{comment}
\usepackage{titlesec}
\usepackage{csquotes}
\usepackage[space]{grffile}
\doublespacing
\usepackage{dcolumn,booktabs}
\interfootnotelinepenalty=10000
\usetikzlibrary{snakes}
\usetikzlibrary{patterns,decorations.pathreplacing}
\usepackage[top=1in, bottom=1in, left=1in, right=1in]{geometry}
\usepackage{bm}
\newcommand{\minus}{\scalebox{0.75}[1.0]{$-$}}
\makeatletter
\renewcommand\@biblabel[1]{}
\makeatother
\usepackage{rotating}
\usepackage{thmtools, thm-restate}
\usepackage{environ}
\usepackage{amsthm}
\usepackage{setspace}
\usepackage{authblk}
\usepackage{hyperref}
\usepackage{pdflscape}
\usepackage{subfig}
\usepackage{natbib}
\usepackage{bbm}
\newcommand*{\QEDA}{\ensuremath{\blacksquare}}
\newcommand*{\QEDB}{\ensuremath{\square}}
\DeclareMathOperator*{\argmax}{argmax}
\newtheorem{prop}{Proposition}
\newtheorem{example}{Example}
\newtheorem{theorem}{Theorem}
\newtheorem{lemma}{Lemma}
\newtheorem{result}{Result}
\newtheorem{assumption}{Assumption}
\newtheorem{corollary}{Corollary}

\hypersetup{colorlinks=true, linkcolor=black, anchorcolor=black, citecolor=black, filecolor=black, menucolor=black, runcolor=black, urlcolor=blue}

\newcolumntype{L}[1]{>{\raggedright\let\newline\\\arraybackslash\hspace{0pt}}m{#1}}
\newcolumntype{C}[1]{>{\centering\let\newline\\\arraybackslash\hspace{0pt}}m{#1}}
\newcolumntype{R}[1]{>{\raggedleft\let\newline\\\arraybackslash\hspace{0pt}}m{#1}}


\hypersetup{colorlinks=true, linkcolor=black, anchorcolor=black, citecolor=black, filecolor=black, menucolor=black, runcolor=black, urlcolor=blue}

\newcommand{\E}{\mathbb{E}}

\usepackage{everypage}

\newlength{\hfoot}
\newlength{\vfoot}
\AddEverypageHook{\ifdim\textwidth=\linewidth\relax
\else\setlength{\hfoot}{-\topmargin}
\addtolength{\hfoot}{-\headheight}
\addtolength{\hfoot}{-\headsep}
\addtolength{\hfoot}{-.5\linewidth}
\ifodd\value{page}\setlength{\vfoot}{\oddsidemargin}
\else\setlength{\vfoot}{\evensidemargin}\fi
\addtolength{\vfoot}{\textheight}
\addtolength{\vfoot}{\footskip}
\raisebox{\hfoot}[0pt][0pt]{\rlap{\hspace{\vfoot}\rotatebox[origin=cB]{90}{\thepage}}}\fi}


\makeatletter
\setlength{\@fptop}{0pt}
\makeatother

\title{\singlespacing{\textbf{What's in a Debt? \\ Rating Agency Methodologies and Firms' Financing and Investment Decisions}}\thanks{We thank Manuel Adelino, Andres Almazan, Aydo\u{g}an Alt{\i}, Jaime Cohen, Jonathan Cohn, Adolfo De Motta, Ben Froehlich, Willie Fuchs, John Hund, Sam Kruger, Jangwoo Lee, Kai Lu, Gregor Matvos, Aaron Pancost, Amiyatosh Purnanandam, Adriano Rampini, Garrett Schaller, Avi Schiff, Michael Sockin, Jason Sturgess, and Sheridan Titman and participants of the MFA, EFA, SFS Cavalcade, UNC/Duke Corporate Finance Conference, WFA, the SEC, McGill and the University of Texas at Austin for the helpful comments and discussions. A previous version of this paper circulated under the title ``Rating Agencies and Leverage Targets''.}}


\author{Cesare Fracassi\footnote{UT Austin, McCombs School of Business. Email: cesare.fracassi@mccombs.utexas.edu.} \hspace{10mm} Gregory Weitzner\footnote{McGill University, Desautels Faculty of Management. Email: gregory.weitzner@mcgill.ca.}}


\date{April 2026}


\begin{document}
\maketitle
\thispagestyle{empty}


\begin{abstract}
\begin{small}
\begin{singlespace}
\noindent In July 2013, Moody's unexpectedly increased the amount of equity credit speculative-grade firms receive for preferred stock from $50\%$ to $100\%$. Firms affected by the rule change were suddenly considered less levered by Moody's, even though their balance sheets did not change. These firms responded by issuing debt to restore the original leverage ratio as defined by Moody's and growing their assets. The rule change transferred value from debt to equity holders and led to an increase in preferred stock issuance. How rating agencies assess risk thus has a significant causal impact on firms' financing, investment, and security design decisions.


\end{singlespace}
\end{small}
\end{abstract}

\clearpage
\newpage

\pagenumbering{arabic}


\begin{center}
    \textbf{What's in a Debt? \\ Rating Agency Methodologies and Firms' Financing and Investment Decisions}

\end{center}

\vspace{2cm}

\noindent In July 2013, Moody's unexpectedly increased the amount of equity credit speculative-grade firms receive for preferred stock from $50\%$ to $100\%$. Firms affected by the rule change were suddenly considered less levered by Moody's, even though their balance sheets did not change. These firms responded by issuing debt to restore the original leverage ratio as defined by Moody's and growing their assets. The rule change transferred value from debt to equity holders and led to an increase in preferred stock issuance. How rating agencies assess risk thus has a significant causal impact on firms' financing, investment, and security design decisions.

\newpage


\section{Introduction}

Assessing a firm's creditworthiness is a complex task that routinely challenges financial market participants. Because investors and regulators rely heavily on credit ratings to measure credit risk, rating agencies can exert substantial influence on firms' financing decisions. This paper shows that the methodologies and standards used by rating agencies to measure creditworthiness have a causal, pervasive effect on firms' financing and investment decisions, asset prices, and the design of fixed income securities.

The majority of metrics used to measure the ability of a firm to meet its financial obligations (e.g., financial leverage, or debt/EBITDA ratio) include variables such as debt and capital. While such metrics are usually not controversial, how to define debt, and how to empirically measure it, is a source of debate (e.g., see \citealp{welch2011two}). Rating agencies routinely adjust the debt of companies by also considering liabilities that are not defined as debt according to GAAP standards, such as defined benefit pension plans, operating leases, inventory on LIFO cost basis, and hybrid securities.\footnote{\cite{eisfeldt2008leasing} and \cite{rampini2013collateral} discuss the discrepancies between the accounting and economic treatment of leases.} In particular, the classification of hybrids, which carry both debt- and equity-like characteristics, is complex and subjective. Rating agency debt adjustments can materially affect leverage ratios, influencing a firm's credit risk assessment. Further complicating the estimation of credit risk, rating agencies use different methodologies to adjust debt for non-debt liabilities.

Given the subjective nature of these debt adjustments, we investigate how credit risk methodologies used by rating agencies influence firms' financial decisions and shape the fixed income market. We exploit an unexpected change to Moody's methodology on the equity treatment of preferred stock as an exogenous shock to firms' perceived credit risk. Specifically, on July 31, 2013, Moody's increased the equity credit that speculative-grade, or junk, non-financial corporate issuers receive for preferred stock from  $50\%$ to $100\%$. Before the rule change, if a speculative-grade company had \$100m in preferred stock, \$50m counted as debt, and \$50m as equity. After the rule change, all \$100m counted as equity.\footnote{Moody's also decreased the amount of equity credit junior subordinated debt receives, but only five speculative-grade, non-financial firms had junior subordinated debt in their capital structures according to Capital IQ. These firms did not have preferred stock, so they would not affect the main treatment group we focus on.}  Speculative-grade firms with preferred stock were suddenly deemed less levered by Moody's, and thus more creditworthy, even though their balance sheets did not change. Among speculative-grade firms with preferred stock at the time of the rule change, preferred stock made up an average of $9.6\%$ of their book capital, leading to a sudden exogenous $4.8$pp decline in Moody's leverage from an average initial level of $61.9$\%.


To analyze the causal effect of the rule change, we construct a panel of firms affected by the rule change, i.e., with preferred stock in their balance sheet and rated speculative grade by Moody's as of July 2013. At the time, 475 US non-financial public firms were rated speculative (Ba1 or below) by Moody's, of which 44 (9.3\%) had preferred stock. To estimate what the financing decisions of these firms would have been absent the rule change, we choose a set of companies unaffected by the rule change, i.e., speculative-rated firms with no preferred stock, but with similar characteristics and in the same industries as the treated firms. After matching, the treated and matched untreated firms are statistically indistinguishable based on balance sheet, income statement, ratings and leverage variables at the time of the rule change.

First, we compare the debt levels and the leverage ratios of treated and matched untreated firms around the rule change. We find that after the change in methodology, treated firms increased their debt levels by $22\%$, with long-term debt driving the increase. As a consequence, firms' debt to capital ratios increase by $3.1$pp. We find a sharp increase in leverage during the first quarter after the rule change, with a small upward drift afterwards.


Second, after the rule change exogenously lowered their Moody's leverage ratios, firms issued debt over the following two years until their Moody's leverage returned to its level prior to the rule change. Despite large increases in debt, we find no statistically significant change in the Moody's ratings of affected firms.


What is the source of the increase in debt capacity? One possible explanation is that the rule change revealed new information about the creditworthiness of speculative-grade firms with preferred stock, deeming them less risky than previously thought, and thus increasing their debt capacity. Alternatively, the rule change may not have revealed information about risk; rather, investors may have wrongly interpreted an increase in ratings as new information because they rely on credit rating agencies to avoid costly information acquisition (e.g., \cite{ramakrishnan1984information} and \cite{millon1985moral}). Finally, the rule change could have relaxed firms' constraints to raise new debt by creating slack in rating-based contracts. For example, derivative contracts often require firms to post additional collateral if they are downgraded. Moreover, bank loans and bonds are frequently structured with increases in required coupon payments, puts, or collateral calls that are triggered once a firm drops below a particular rating. Ratings triggers, like other covenants, can be used to disincentivize shareholders from taking actions that destroy firm value, such as risk-shifting. Importantly, these three mechanisms are not mutually exclusive, and they all could drive the response to changes in rating methodology.

According to all three mechanisms, a higher rating leads to increased debt capacity. We thus examine the asset side of the balance sheet as well as stock and bond price responses around the announcement of the rule change. We find that firms used the increase in debt capacity to expand their balance sheet through asset growth and investment, rather than simply recapitalizing. The additional debt issuance led to an 8\% increase in property, plant and equipment (PP\&E) and assets. We also find that the stock prices of affected firms increased by an average of 2.8\% in a short window around the announcement of the rule change, while their credit spreads widened by 16bps or 3.5\%. These results are inconsistent with an information-based channel, as we would expect a recapitalization and an increase or no change in bond prices as the perceived creditworthiness of firms increases. The increase in investment and the wealth transfer from bondholders to equityholders suggest that constrained firms were able to expand their investment while making their liabilities riskier, without leading to contractual increases in their cost of capital from being downgraded. While we expect firms to adapt provisions for new liabilities to the new rating methodologies, the existing liabilities are not automatically readjusted, thus firms increase their investment at the expense of debtholders once their existing rating triggers are relaxed. Moreover, the fact that firms issue more debt after existing bond yields rise is consistent with bond markets anticipating future increases in debt issuance due to the relaxation of their covenant constraints. The market reaction also shows that Moody's rule change was indeed unexpected.

To further validate the covenant channel, we develop cross-sectional tests based on \cite{green1984investment}, who shows that convertible debt reduces the incentives for firms to engage in risk-shifting. Intuitively, the more of a firm's debt that is convertible into equity, the lower the benefits of risk-shifting because convertible debtholders receive a portion of these benefits. Hence, we test whether firms respond less to the rule change the larger the fraction of their debt that is convertible debt. Consistent with our hypothesis, we find that firms with more convertible debt issue less debt, increase their leverage and invest less following the rule change. In summary, the results are consistent with the view that Moody's methodology change relaxed rating-based covenants and firms' issuance, leverage and investment responses were partially driven by the value they could extract from debtholders. Nonetheless, while we cannot fully rule out other channels, our tests suggest that the covenant channel is the dominant one.


Finally, the rule change not only influenced the choice of debt and investment levels but also ultimately affected the type of securities issued by firms. The change in Moody's methodology made preferred stock more attractive for speculative-grade firms, as it became treated as pure equity. We find that firms rated speculative-grade only by Moody's nearly tripled their preferred stock levels relative to firms rated speculative-grade by both Moody's and S\&P, or only by S\&P, following the rule change. As a placebo test, we perform the same analysis on investment-grade firms and find no change in preferred stock issuance.

We conduct a range of robustness tests to confirm the internal validity of our results. First, the trend in outcome for the period before the rule change is similar between the treated and control groups, i.e., parallel trends hold. Second, our results remain if we use all speculative-grade firms without preferred stock as a counterfactual, rather than matching on firm characteristics. We also find similar results when using investment-grade firms with preferred stock as a counterfactual control group. We also conduct several placebo tests: we pretend the rule change occurred in July of other years besides 2013 and find no statistically significant change in debt levels outside of 2013. We also find no results when we instead use convertible debt, another hybrid security with features similar to preferred stock, as the treatment variable. We also find no effect when comparing firms with preferred stock to those without, among unrated firms rather than speculative-rated firms.\footnote{Like preferred stock, convertible bonds may be used to avoid the financial distress associated with pure debt and mitigate adverse selection costs (e.g., \citealp{stein1992convertible}).}


It is important to highlight the limitations of the study. First, we cannot fully generalize the findings beyond the firms specifically affected by the rule change. These firms have high levels of debt and, thus, might be more sensitive to the influence of rating agencies than the average public firm. However, we show that the treated firms in our sample have levels of financial constraints similar to those of all other speculative-rated firms. Second, our channel relies on the fact that firms are financially constrained. Unconstrained firms would not be affected by rating methodology changes. Third, preferred stock is not allocated to firms randomly; thus, our difference-in-differences strategy relies on the following conditional independence assumption: after controlling for observable characteristics, a firm's choice of having preferred stock in its balance sheet in July 2013 is orthogonal to its financing and investment decisions made afterward. While this identifying assumption cannot be directly tested, we show that treated and matched control firms have similar levels of observable characteristics, exhibit parallel trends in the outcome before the rule change, and that treated firms show no increase in debt relative to control firms in any year other than 2013. Furthermore, the issuance of preferred stock occurs on average 9.5 years before the rule change, and 70\% of the preferred stock is non-callable. Thus, preferred stock is a perpetual instrument that, in many cases, was on the balance sheet for many years before the rule change. Finally, asset prices reacted in a narrow window around the rule change, suggesting that neither credit nor equity markets anticipated it. Moreover, we find no evidence of other major announcements or macro-events occurring around the time of the rule change.

The paper contributes to various strands of finance literature. While several papers have studied the effect of ratings on firm financing decisions (among many, \citealp{graham2001theory}, \citealp{kisgen2006credit}, \citealp{sufi2007real}, \citealp{kisgen2009firms}, \citealp{hovakimian2009credit}, \citealp{begley2013real}, \citealp{adelino2017economic} and \citealp{almeida2017real}), our paper shows that rating agencies also influence firms through the way in which they define credit risk. We also provide evidence of the channel through which ratings influence firms' decisions. Another paper analyzing changes in rating methodologies is \cite{kisgen2019impact}, which uses Moody's 2006 changes in the treatment of underfunded pensions and operating leases. Our paper differs significantly in four main ways. First, the change in rating methodologies applied only to specific firms, clearly defining treatment and control groups. Second, we expand the scope of the study by examining stock prices and bond spreads, providing market-based evidence that the rule change had real valuation consequences. Third, we show that firms actively redesigned their hybrid securities in response to the reclassification. Finally, we find evidence of a risk-shifting mechanism, documenting a potential channel through which changes in rating methodology affect firm behavior.

We also contribute to the debate as to the speed at which firms respond to slack in their debt capacity  (e.g. \citealp{fama2002testing}, \citealp{baker2002market}, \citealp{welch2004capital}, \citealp{leary2005firms}, \citealp{flannery2006partial}, \citealp{alti2006persistent} and \citealp{kayhan2007firms}). One challenge in the literature is that leverage targets are endogenous and can be time-varying. It is also difficult to distinguish managing leverage to a target from alternate motives (\citealp{iliev2010reconciling} and \citealp{graham2011review}). Our paper uses an identification strategy in which leverage changes, but firms' balance sheets do not.\footnote{Or equivalently, GAAP leverage targets change.} We find that it takes two years for firms to return to the original leverage ratio as defined by the rating agencies, with a $30\%$ adjustment occurring in the first quarter, a faster speed of adjustment than the estimates found in the existing literature, ranging from $0\%$ to $39\%$ per year.

Recent papers also question the objectivity of rating agencies in assigning credit ratings (e.g. \citealp{griffin2012did}, \citealp{fracassi2016does}, \citealp{cornaggia2016revolving}, \cite{kempf2018partisan}, and \citealp{wang2021rating}). In our setting, the treatment of preferred stock appears subjective, given that S\&P and Moody's use vastly different methodologies. Furthermore, there is evidence that standards seem to change over time. \cite{alp2013structural} finds that investment-grade ratings tightened while speculative-grade ratings loosened over the period of 1985-2002. Similarly, \cite{baghai2014have} find that credit rating agencies have become more conservative over time. Our analysis shows that subjectivity in credit methodologies influences firms' financing, investment and security design decisions.


Finally, we contribute to the empirical literature on shareholder/debtholder conflicts, risk-shifting, and covenants (\citealp{eisdorfer2008empirical}, \citealp{gilje2016firms}, \citealp{rauh2008risk}, and \citealp{gormley2011growing}). In our paper, we find evidence consistent with the rule change relaxing rating-based triggers for a subset of firms, allowing them to engage in risky investment at the expense of existing debtholders. Furthermore, the response is stronger among firms with a higher incentive to risk-shift. The majority of the existing empirical literature on covenants analyzes ex-post violation of covenants, while we arguably identify an interim relaxation of covenants.\footnote{For example \cite{chava2008does}, \cite{roberts2009renegotiation} and \cite{nini2009creditor}.} \cite{matvos2013estimating} and \cite{green2018corporate} estimate structural models and find large ex-ante benefits of covenants. \cite{admati2018leverage} show theoretically that leverage ratios are inherently unstable in the absence of covenants, which is consistent with the increase in leverage we observe once rating triggers are relaxed. Our analysis suggests that tying contracts to ratings may be an imperfect tool for mitigating conflicts between debtholders and shareholders due to the sometimes arbitrary nature of rating agency methodologies.

\section{Background on Rating Agency Debt Adjustments}

The amount of debt of a firm may not reflect its underlying credit risk (\citealp{graham2011review}). Items such as operating leases, unfunded pensions, capitalized interest, and hybrid securities may not appear as debt on a firm's balance sheet but act as ``debt-like'' obligations. Rating agencies recognize this discrepancy and use debt adjustments to better assess a firm's creditworthiness. Such debt adjustments influence many of their own internal credit risk metrics, such as leverage and debt-to-EBITDA ratios.


Hybrid securities, such as preferred stock and junior subordinated debt, are particularly difficult to evaluate. Under GAAP accounting, preferred stock generally appears on the balance sheet as equity (\citealp{pwc}). However, rating agencies consider preferred stock a hybrid due to the recurring fixed payment, which is similar to that of a debt instrument, and its seniority to common equity. Rating agencies use hybrid basket treatment tables to assess the amount of equity credit a specific preferred stock instrument receives. For example, a preferred stock with a finite maturity is given less equity credit than perpetual preferred stock by both Moody's and S\&P.


There are several reasons why firms may issue preferred stock. First, if the distress costs from pure debt are high, firms can issue preferred stock when adverse selection is high, making equity costly to issue, or if managers believe their stock is undervalued. Second, preferred stock allows firms to issue permanent capital while retaining control, as most preferred equity does not carry voting rights. Third, preferred stock may also be issued for clientele reasons as preferred dividends face lower tax rates for investors than interest payments (\citealp{RAVID2007198}). \cite{heinkel1990role} show that preferred issuance can increase total debt capacity. \cite{nance1993determinants} argue that preferred stock can act as a substitute for hedging by reducing the probability of financial distress. Empirically, \cite{RAVID2007198} analyze the determinants of preferred stock usage in an older sample period and finds some evidence that poorly performing companies are more likely to issue preferred stock, which is consistent with the financial distress and asymmetric information rationales. In Table \ref{table_prefdet}, we analyze the determinants of preferred stock in our more recent sample and find fairly similar results.


\subsection{Rule Change}
Before July 2013, for both Moody's and S\&P, \$100 of perpetual preferred stock counted as \$50 of equity, and the other \$50 was counted as debt for non-financial issuers of all ratings. On July 31st, 2013 Moody's issued a report altering its leverage adjustments for hybrid securities, stating the following:

\begin{displayquote}
\begin{small}
\begin{singlespace}
``Relative to investment-grade nonfinancial companies, speculative-grade nonfinancial companies are materially closer to default, have shorter dated and more complex capital structures, as well as debt with more covenants. Additionally, speculative-grade nonfinancial companies often opt to cease cash distributions associated with preferred stock and other hybrid instruments because such actions are contractually allowable without triggering a debt default...

...Given these characteristics, our approach to assessing the debt and equity characteristics of hybrid instruments of speculative-grade companies with a corporate family rating (CFR) or senior unsecured rating of Ba1 and below is to closely follow the legal treatment we expect these instruments would receive in a bankruptcy scenario. Instruments with a debt claim as set out herein, other than shareholder loans meeting defined criteria, receive 100$\%$ debt treatment while equity instruments, such as preferred stock, with no such debt claim receive 100$\%$ equity treatment'' \cite{moodydebt}.
\end{singlespace}
\end{small}
\end{displayquote}

The rule change thus increased the amount of equity credit that Moody's provides perpetual preferred stock for speculative-grade corporate issuers from $50\%$ to $100\%$, e.g., \$100 of perpetual preferred stock would now count as \$100 of equity. At the time, Moody's rated 475 US, non-financial, public firms as speculative (Ba1 or below), of which 44 (9.3\%) had perpetual preferred stock. On average, speculative-rated firms with preferred stock had \$2.7bn in debt, \$395mil in preferred stock, and \$5.2bn in capital. Before the rule change, half of the \$395mil in preferred stock was counted as debt by Moody's, and none after the rule change. Suddenly, speculative-grade firms with preferred stock were deemed less levered by Moody's, even though their balance sheet did not change. The average implied change in leverage as measured by Moody's declined by 4.8pp from $61.9\%$ to $57.1\%$, which corresponds to more than a notch upgrade in credit ratings.\footnote{We estimate that one notch corresponds to approximately 5pp in leverage. See Table \ref{table_ratbook} in the Online Appendix for more details.} Figure \ref{fig_histo} shows the distribution of preferred stock over capital among speculative-grade firms affected by the rule change.


The rule change was permanent, as Moody's affirmed its new treatment in March 2015. Moody's did not change the equity credit given to preferred stock for investment-grade issuers and S\&P maintained the original equity credit for both investment-grade and speculative-grade issuers. The rule change thus provides several natural control groups to test in the empirical analysis. Moody's also reduced the equity credit junior subordinated debt receives for speculative-grade firms from $25\%$ to $0\%$. Following the rule change, ArcelorMittal reportedly called an existing junior subordinated note because it no longer received equity credit \cite{ipe}.\footnote{We do not analyze junior subordinated debt because speculative-grade firms rarely use it. Only five issuers with junior subordinated debt were in Capital IQ at the time of the rule change, and none of these firms had preferred stock in their capital structure, which could contaminate our treatment group.}


Investors appear to be aware that rating methodologies can change over time and acknowledge the potential for disruption in fixed income markets. ``The behaviour of rating agencies is one of the biggest risks,” said Anne Velot, head of Euro credit at AXA Investment Managers. Robert Emes, a credit analyst at JPMAM, said: “While not all changes will cause an issuer to redeem early, the agencies will constantly try to reform their methodologies to make them more indicative of the risks that they see, so I can’t rule out that this might happen.”  While investors might be aware of the possibility of changes in rating methodologies, news surrounding the event described the rule change as a surprise to the market. For instance, a market commentator stated: ``Moody’s decision to remove equity credit from all hybrid bonds issued by a company rated Ba1 or lower in August 2013 caused turbulence and was a reminder that rating methodologies can be subject to rapid change'', \citep{euromoney}.\footnote{This specific quote refers to the fact that Moody's removed equity credit for junior subordinated bonds at the same time that it increased it for preferred stock. On November 27$^{th}$ 2015, S\&P also suddenly changed the equity credit for a small group of European hybrids. See \cite{euromoney}}

Why did Moody's change its treatment for hybrids? The report stated that for investment-grade firms, preferred stock dividends are more often uninterrupted, while riskier, speculative-grade firms that are closer to default often cease preferred dividend payments. Thus, they deemed preferred stock more debt-like for investment-grade firms and more equity-like for speculative-grade firms. Since preferred stock dividends cannot trigger default and have no debt claim in bankruptcy, Moody's decided to treat preferred stock as equity for speculative-grade firms.

\subsection{Methodology to Measure Leverage Ratio}\label{sec_ral}


GAAP leverage, the most commonly used leverage ratio by academics and practitioners, is defined as follows:


\begin{align}\label{lev_gaap}
Lev_{GAAP} =  \frac{D}{D+E+P},
\end{align}
where $D$ denotes the face value of debt, $E$  the book value of the firm's equity besides preferred stock and $P$ the principal of the firm's preferred stock. In defining leverage, we follow \cite{welch2011two} and avoid the debt-to-asset ratio; however, the main results hold with this measure and are displayed in the Online Appendix.

As discussed above, credit rating agencies make debt adjustments to reflect debt-like liabilities in firms' balance sheets. In this study, we only analyze debt adjustments related to preferred stock, as these adjustments were suddenly changed by Moody's. Moody's leverage is thus defined as:

\begin{align}\label{RAlev1}
Lev_{Moody's} =
\begin{cases}
 \frac{D +0.5P}{D+E+P} & \textnormal{ before July 31st, 2013 } \\
\frac{D}{D+E+P}  & \textnormal{ after July 31st, 2013.}
\end{cases}
\end{align}


\section{Hypothesis Development and Research Design}

\subsection{Rating Methodologies and Firm Decisions} \label{sec_firmdec}

In our view, three main channels can explain how rating methodologies affect a firm's financing and investment decisions.

Many financial contracts, such as bank loans, derivative contracts, leases, and trade agreements, include terms contingent on the creditworthiness of the parties signing the contract. Given that credit spreads are for the most part not available or sufficiently liquid, contracts often rely on the credit rating of a firm.\footnote{CDS are generally more liquid but are only traded on a small universe of issuers.} Covenants can be used to prevent shareholders from taking actions, such as risk-shifting, that benefit shareholders, but decrease firm value when risky debt is in place (\citealp{smith1979financial}).\footnote{See also \cite{bhanot2006should} for a theoretical motivation of rating triggers when there is a risk-shifting problem.}  Counterparties can use rating triggers to prevent firms from taking on risky investment by causing a sudden increase in the security's coupon, a collateral call, or debt acceleration when the firm is downgraded.\footnote{
A rating trigger may be more useful to prevent risk-shifting than a financial leverage covenant, such as debt/EBITDA, as a firm could easily increase its average future risk without affecting its debt/EBITDA ratio, while credit ratings reflect the expected future risk of the firm's assets. \cite{parlour2018contracting} rationalize the use of ratings for contracting purposes in an investment management context when contracts are incomplete.}  In 2002, Moody's conducted a survey of US corporate issuers and found that 87.5\% of respondents had ratings triggers with an average of over four triggers per company (\citealp{moodytrigger2}). According to this view, we would expect firms to respond to the rule change by increasing leverage and investment at the expense of existing debtholders and counterparties, causing an increase in stock prices and a decrease in bond prices around the rule change.\footnote{Firms may also simply issue more debt and use the proceeds to return capital to shareholders, diluting existing liabilities. However, leverage covenants would seem to be a better tool for preventing this.} We also expect that firms' Moody's ratings remain the same, and that firms that are not at risk of being downgraded by S\&P will respond more, consistent with firms targeting a minimum credit rating to avoid rating-based triggers.


The second channel through which rating methodologies can affect firms is new information revealed by rating agencies. Many investors rely on ratings to assess credit risk. \cite{tang2009information} finds that after Moody's transition from broad rating categories to narrower notches, credit spreads responded depending on the direction in which the rating moved. Firms and investors may view Moody's announcement as a real reduction in credit risk for the affected firms because Moody's revealed new information about the riskiness of preferred stock for speculative-grade firms. Suppose a firm is at its optimal leverage under the trade-off theory. After the announcement, the firm and markets learn that its expected cost of distress is lower; thus, the firm issues more debt to equate the marginal cost of distress with the marginal benefit of tax shields. Since the new information is about the riskiness of the firm's liabilities, not its marginal investment opportunities, we would expect a recapitalization rather than an increase in investment. In terms of asset prices, we would expect bond prices to either increase due to the new information or remain flat if investors anticipate that firms will issue more debt afterward to offset the decrease in credit risk. We would also expect equity prices to increase due to new information about the riskiness of the firm's liabilities and its ability to issue more debt to take advantage of higher tax shields.

The third way rating methodologies could influence firms' behavior is through limited attention. Firms and investors may not be aware of the rule change and may merely follow the firm's credit rating to assess credit risk. The reliance of stakeholders and investors on ratings could reduce costly individual information acquisition (e.g., \cite{ramakrishnan1984information} and \cite{millon1985moral}). Therefore, investors and firms would need to see changes in ratings to spur changes in firm behavior and asset prices. Under this channel, we would expect firms to increase their debt and investment following an upgrade because they would wrongly believe they are less risky. Similarly, if ratings immediately rise, we would predict the same asset price responses as in the rational information channel: stock prices would increase, and bond prices would either increase or remain flat. However, under this channel, we would only expect to see a response if firms' ratings actually change after the rule change.


Finally, there may also be regulatory frictions as explored in \cite{bruno2015does}, \cite{kisgen2010regulations}, \cite{almeida2017real} and \cite{bernstein2017costs}. For instance, certain institutional investors may be required to hold investment-grade bonds.\footnote{Very few of the affected firms in our sample are even near the investment-grade boundary.} In addition, capital requirements often depend on ratings for financial firms.\footnote{Hybrids are popular securities among financial firms because of the regulatory capital they receive from them. However, the rule change we analyze in this paper only applies to non-financial firms.} However, these channels are less likely among speculative-grade, non-financial firms. \cite{almeida2017real} exploits corporate downgrades stemming from rating agency policies requiring corporate ratings to remain below the sovereign rating of their country of domicile. Given that the affected firms in our natural experiment are all below investment-grade and domiciled in the US, this policy does not apply.

Providing evidence on the channels through which rating methodologies affect firms’ financing and investment decisions is crucial for both understanding the influence of rating agencies and assessing the broader applicability of our results. For instance, if the primary channel operates through covenants, we would expect our findings to extend to other firms that face similar covenant constraints, even for those not based on credit ratings.


\subsection{Experiment Design} \label{sec_experiment}


Moody's change in the treatment of preferred stock only applied to a subset of US public firms, allowing us to compare firms affected by the rule change with similar firms that were not affected. First, we select firms that were affected by the rule change. Using Compustat Fundamentals Quarterly, Thomson Eikon, CRSP, and Intercontinental Exchange (ICE), we collect financial data for all US public non-financial firms with a speculative-grade rating from Moody's (Ba1 or lower) as of July 31st, 2013, yielding 475 firms.\footnote{We exclude financial firms and non-operating establishments (SIC codes with 6000 - 6999 and 9995) because financial firms are not affected by the rule change; there is one firm with preferred stock in non-operating establishments, but it is classified as a financial institution by Moody's; thus, for consistency, we remove all non-operating establishments.} Among these firms, we then proceed to select those that have a positive amount of preferred stock (Compustat item prstkq) on their balance sheet in the last quarter prior to July 2013. We hand-check these firms and remove any that have only mandatory convertible preferred, trust preferred, preferred with any maturity or puts, or no preferred at all, as these types of preferred are not treated as perpetual preferred stock by Moody's and were not subject to any changes in treatment. We search through reports by Moody's to see if there is mention of Moody's treatment of particular securities.\footnote{For instance, according to Compustat iGate had a positive amount of preferred stock before the rule change; however, the preferred instrument had a put, thus it was treated as debt by Moody's.``Moody's views the $\$350$ million of preferred stock issued to the financial sponsor as debt like due to the six year maturity at which time the holders may require iGate to redeem its shares for cash equal to the accrued liquidation preference" (\citealp{moodyassigns}).} We also collect the proper principal value from SEC 10-Q and 10-K filings, as Compustat often has the incorrect face value. Overall, we find 44 firms that are both rated speculative-grade by Moody's and have perpetual preferred stock on their balance sheet. These firms constitute our treatment group.

Second, we identify firms that are as similar as possible to the treated firms, but were not affected by the rule change. We draw from the pool of speculative-grade rated firms without preferred stock on their balance sheet to find our main counterfactual. Table \ref{table_sum} shows that firms affected by the rule change (treated firms) have lower market-to-book, lower profitability and worse ratings relative to the rest of the speculative-rated firms. We thus use a matching strategy to select the set of control firms used as a counterfactual. For each treated firm, we find four speculative-grade firms in the same Fama-French 49 industry that are as similar as possible in terms of Moody's ratings, growth in leverage, and market-to-book at the time of the rule change.\footnote{We also construct an alternative matched sample in which each treated firm is paired with its single closest control firm, without replacement. The results from this alternative matching procedure are reported in Online Appendix Tables \ref{table_debt_alt}-\ref{table_bs_alt}. They are remarkably similar to our main findings, both economically and statistically, reinforcing the robustness of our conclusions.} As is standard in the literature, we use the Mahalanobis metric as a distance measure, i.e., the inverse of the sample variance/covariance matrix of the three continuous non-exact matching variables. We also control for profitability, tangibility, market-to-book and log of sales in all specifications, as these are the most common controls in the capital structure literature (e.g. \citealp{rajan1995we} and \citealp{deangelo2015stable}).  The precise definition of all variables used in the paper is shown in Table \ref{table_vardef}. We winsorize the control variables at the $1\%$ and $99\%$ levels to make sure that outliers do not influence our results.

The rule change provided companies with financial slack, increasing their debt capacity. A possible concern is that firms with preferred stock might be systematically more financially constrained than the average firm. If financial constraint is the primary channel driving our results, the findings might not generalize to the broader population of speculative-rated firms. To address this potential threat to external validity, we compare treated, untreated control firms, and all other speculative-rated firms using three measures of financial constraints: Kaplan and Zingales (KZ) (\cite{Lamontetal2001}), Whited and Wu (WW) (\cite{WhitedWu2006}), and Size - Age (\cite{HadlockPierce2010}). Online Appendix Table A.2 and Figure A.1 suggest that treated and non-treated firms have similar levels of financial constraints.


Our identification strategy relies on the parallel trend assumption that after we match and control for observable characteristics, the unobserved time-series trends in the financing and investment decisions of firms are orthogonal to whether these firms had preferred stock in their balance sheet in July 2013. While the main identifying assumption cannot be tested explicitly, three pieces of evidence are consistent with this assumption: (i) As shown in Table \ref{table_sum}, the difference between the treated and the matched untreated firms is not statistically significant not only for the variables used for matching, but also for other observables; (ii) in all regressions, we test for parallel trends before the change in Moody's methodology, and there is no difference in the time-series behavior of treated and matched untreated firms; and (iii) because the decision to issue preferred stock occurred on average 9.5 years before the rule change, it is thus unlikely that unobservable variables that lead firms to issue preferred stock a decade earlier will impact the change in balance sheet only in July 2013, but not before.\footnote{These statistics come from hand-collected issuance data on the preferred stock of the treated firms.}


An alternative, and in our opinion weaker, approach would be to use investment-grade firms with preferred stock as the control group. While this approach controls for firms' decisions to use preferred stock, it relies on the assumption that a firm's investment-grade status is orthogonal to its financing and investment decisions, an assumption that is harder to justify. Nonetheless, in Tables \ref{table_diff_c} and \ref{table_convertIG}, we show that results are robust to using such an alternative counterfactual.

Finally, after having identified cohorts of treated and matched control firms, our main empirical analysis employs a cohort generalized difference-in-differences strategy. Essentially, we take the difference in outcome $y_{i,c,t}$ for each treated firm $i$ after the rule change relative to before and compare it with the difference in outcome of its matched control firms within the same cohort $c$.
	\begin{equation} \label{eq:single}
	y_{i,c,t}= \beta (d_{i,c} \times p_{t,c}) +  \gamma X_{i,t} + \alpha_{i,c} + \delta_{t,c} + u_{i,c,t}.
	\end{equation}
All regressions are estimated from 4 quarters before the July 2013 event to 8 quarters afterward. We choose the pre-event window to have enough periods to test the pre-event parallel trends. We select the post-event window to allow firms sufficient time to respond to the rule change. It is also consistent with the speed of adjustment tests found in \cite{leary2005firms}. In all tables, we use two types of treatment variables: (i) a discrete dummy treatment variable $d_{i,c}$ equal to one if the firm has a positive amount of preferred stock at the time of the rule change, and (ii) a continuous treatment variable defined as the amount of preferred stock relative to capital that the firm had at the time of the rule change.\footnote{For all firms unaffected by the rule change, this variable equals zero.} $p_{t,c}$ is a dummy variable equal to one if the time period is after the rule change. $X_{i,t}$ are the control variables profitability, market-to-book, tangibility, and log of sales. The coefficient $\beta$ represents the difference-in-differences effect of the rule change on the outcome variable relative to a matched counterfactual. The firm-cohort fixed effect $\alpha_{i,c}$ ensures that we compare changes in outcomes within the same firm.\footnote{To alleviate concerns that the control variables were also affected by the rule change, causing a bias in the coefficients, we also estimate regressions keeping the value of the control variables constant as of the time of treatment throughout the sample period. In untabulated results, we find that the economic magnitude of the point estimates is even larger.} The time-cohort fixed effect $\delta_{t,c}$ ensures that the treatment firm is compared only with the four matched control firms at each point in time. The standard errors are clustered at the firm level to adjust for heteroskedasticity and serial correlations in the error term (\citealp{petersen2009estimating} and \citealp{thompson2011simple}).

We also estimate the impact of the rule change quarter-by-quarter, using the equation below:
	\begin{equation} \label{coefficients}
	y_{i,c,t}=\sum_{k=-4}^{8} \beta_k (d_{i,c} \times \lambda_{t,k,c}) + \gamma X_{i,t} + \alpha_{i,c} + \delta_{t,c} + u_{i,c,t},
	\end{equation}
where $\lambda_{t,k,c}$ is a dummy equal to one if the calendar quarter $t$ is equal to $k$ and zero otherwise, where $k=0$ is the quarter of the rule change. Standard errors are also clustered at the firm level. The quarter-by-quarter regressions help ensure that the evolution of the outcome variable before the event is similar across the two groups. A positive or negative pre-event trend could invalidate the interpretation of the difference-in-differences results. In addition, we can learn how quickly firms respond to the rule change and infer the speed of adjustment to shocks to leverage.  Given the large number of fixed effects and observations, all regressions in the paper are estimated using the fixed point iteration procedure implemented by \citet{Correia2017:HDFE}.

Table \ref{table_sum} shows the summary statistics as of the last quarter prior to the rule change for the three groups of firms in the sample: the treated firms (the ones affected by the rule change), all untreated firms (all the speculative-grade firms which were not affected by the rule change) and the matched firms (the subset of untreated firms which are as close as the treated firms as possible). There are 44 speculative-grade firms with preferred stock at the time of the rule change, which is about 9\% of the total number of speculative-grade firms in the sample. Table \ref{table_list} shows the complete list of treated firms and their Fama-French 49 industry classification. Treated firms belong to a variety of industries, including energy, housing, media, industrial and retail. The characteristics of the treated and matched firms are fairly similar at the time of the treatment: only S\&P ratings are statistically different at the 10\% level; however, if we use a Šidák correction to control for false discoveries, we fail to reject the null hypothesis that the treated and matched control variables are the same across all variables. The average book GAAP leverage of the firms is $57.1\%$, corresponding to a credit rating of B1 using the Moody's rating scale.\footnote{Moody's and S\&P issuer level ratings are translated to a number with 1 being the least likely and 22 being the most likely to default.} Hence, the average rating of firms in our sample is well below investment grade, and only four firms are within two notches of investment grade. This is important, as firms near the investment-grade threshold may have a strong incentive not to issue debt after the rule change, as it could be an opportunity to be upgraded to investment-grade, thereby achieving a lower cost of capital because of regulatory and institutional frictions. The average market capitalization is $\$3.0$bn, slightly smaller than that of the average public firm in the Compustat universe, $\$3.8$bn.


\section{Results}

\subsection{Moody's Rule Change and Firms' Financing Decisions}\label{sec_fin}

Moody's change in the treatment of preferred stock led to several speculative-rated firms being deemed less levered by Moody's, even though their balance sheet did not change. We thus investigate whether firms responded to this rule change by altering their capital structure using the difference-in-differences specifications in Equations (\ref{eq:single}) and (\ref{coefficients}).

First, Table \ref{table_debt} shows that treated firms increased their use of debt after the rule change, relative to a matched group of comparable speculative-grade firms without preferred stock. On average, the total debt held by treated firms increased by over 22\% after the rule change (Column (1)). The effect is greater for firms with a large amount of preferred stock in the balance sheet at the time of the rule change (Column (2)). The effect is statistically significant for long-term debt, and not for short-term debt (Columns (3) to (6)). Thus, firms actively issue bonds and bank debt to increase their leverage ratio. Figure \ref{fig_debt_reg} shows the trend over time of the increase in total debt. Treated firms have comparable levels of debt relative to matched control firms in the year before the rule change, but then the difference rapidly increases in the first two quarters after the rule change, before plateauing.

Table \ref{table_placebo} shows the results of a placebo test, where we rerun the regression specification in Equation (\ref{eq:single}) for each calendar year from 2008 to 2015, using total debt as the dependent variable. To avoid overlapping periods, we restrict the sample period to a single calendar year. The ``After July'' variable is equal to one if the calendar quarter ends after July, and zero otherwise. The table shows that treated firms increase their debt significantly relative to matched control firms only in 2013, which is the year of the rule change. These results show that the evidence presented in Table \ref{table_debt} is not driven by any seasonal pattern, and it does not occur in any other year except the year of the rule change. We also show that the broader equity and treasury markets were virtually unchanged around the announcement of the rule change, making it unlikely our results can be explained by other simultaneous large events (see Online Appendix Table \ref{table_carnews}).

Second, we check whether the rule change led to an increase in GAAP leverage $Lev_{GAAP}$ (Equation (\ref{lev_gaap})). In Table \ref{table_lev}, we find that GAAP leverage increases after the rule change. Column (1) shows that the average book GAAP leverage of speculative-grade firms with preferred stock increases by 3.1pp (from 57.1\% to 60.2\% for the average firm) relative to matched speculative-grade firms without preferred stock. Assuming the firm's book equity is unchanged, the increase in leverage is consistent with the 22\% increase in debt.\footnote{A 22\% increase in debt corresponds to a 4.8pp increase in debt to capital, assuming the firm does not adjust its book equity, which is within the confidence interval of the point estimate for leverage.} The results show that the exogenous rule change led to a significant change in the financing and capital structure decisions of affected firms.


Third, we study the evolution of leverage using Moody's definition $Lev_{Moody's}$ (Equation (\ref{RAlev1})). At the time of the rule change, we should observe a sharp mechanical drop in leverage: given that preferred stock receives 100\% equity credit after the rule, if companies did not respond to the rule change, the Moody's leverage should drop by half of the fraction of capital in preferred stock (i.e the $\beta$ coefficient right after the change should be -0.5). If firms responded immediately by issuing enough debt to offset the drop, Moody's leverage should not change over time. Figure \ref{fig_bl_reg} shows that $Lev_{Moody's}$ drops after the rule change, but not by the full amount triggered by the rule change: firms respond to the rule change quickly, offsetting 30\% of the newly available debt capacity within the first two months. After the first quarter, a slow uptrend emerges over the subsequent seven quarters, until the drop in Moody's leverage becomes statistically insignificant after 18 months and almost completely disappears two years after the rule change. This evidence suggests that firms use up the debt capacity created by changes in rating methodologies, and that they respond quickly to exogenous shocks to their rating-agency leverage. The $30\%$ initial adjustment in the first quarter after the rule change is faster than the existing findings in the speed-of-adjustment literature.\footnote{\cite{graham2011review} synthesizes the findings of \cite{fama2002testing}, \cite{kayhan2007firms}, \cite{flannery2006partial}, \cite{lemmon2008back}, \cite{huang2009testing}, \cite{elsas2015dynamic}, and \cite{iliev2010reconciling} who use different methods to measure the speed-of-adjustment.}


The evidence of an increase in the leverage of treated firms after the rule change supports both the information-based channels and the covenant-based channel. Next, we explore the impact of the rule change on the asset of the balance sheet.


\subsection{Moody's Rule Change and Firms' Investment Decisions}

After Moody's changed its methodology for preferred stock, firms responded by increasing their debt levels and leverage. The additional capital raised could be used to either shrink their equity base through dividends and share buybacks or to increase their balance sheet. We thus test whether the rule change affects the real asset side of the balance sheet. In Table \ref{table_bs}, we use the natural logarithm of capex, property, plant and equipment (PP\&E), and total assets as outcome variables.\footnote{PP\&E is the only variable for which we use the log(1 + x) transformation. In Online Appendix Table \ref{table_bs_sine}, we show that the results are even stronger when we apply the inverse hyperbolic sine transformation (\cite{cohn2022count}).} Column (1) shows that capital expenditures increase by 10\% for treated firms relative to matched control firms; however, the estimate is not quite statistically significant. The increase in capex leads to a statistically significant 8\% increase in PP\&E (Column (3)) and assets (Column (5)). The results are similar in sign and statistical significance when we use a continuous treatment variable rather than a treatment dummy (Columns (2), (4), and (6)). The 8\% increase in assets is consistent with the 22\% increase in debt assuming firms do not use the proceeds to return capital.\footnote{The average debt to assets ratio is 41.7\%, which translates to a 9.1\% increase in assets, which is within the confidence interval of the point estimate for assets (Column (1) of Table \ref{table_bs}).} In untabulated results, we find no increase in share repurchases or dividends, which is also consistent with firms using the proceeds from the debt issuance to increase their investment. Consistent with the time-series trend in debt and leverage levels, Figure \ref{fig_bs_reg} shows that the increase in the balance sheet occurs in the first year, and then it plateaus. Firms appear to be constrained by their ratings, and providing them with extra debt capacity leads to more investments and a larger balance sheet.


Finally, we run a battery of tests to ensure that the results are robust to alternative research design choices. First, we run a placebo test in which we treat speculative-rated firms with convertible debt at the time of the rule change as treated. Like preferred stock, convertible debt is a hybrid security that provides firms with greater financial flexibility.\footnote{See Tables \ref{table_prefdet} and \ref{table_convdet} for an analysis of the determinants of preferred stock and convertible debt.} However, the rule change only affected preferred stock. If unobservables, rather than the rule change, were to drive the change in the balance sheets of firms with preferred stock after July 2013, we might expect similar behavior for firms with convertible debt. We thus create a variable \textit{Convert}, which is the ratio of convertible debt to total debt at the time of the rule change. The odd columns of Table \ref{table_combinedplacebo} show that speculative-rated firms with convertible debt do not behave differently relative to firms without convertible debt.

Second, we repeat the main analyses, but considering only unrated firms, rather than firms that are speculative-rated by Moody's. Unrated firms share characteristics with speculative-rated firms, as both are likely to be financially constrained, unprofitable, and small. If firms with preferred stock behaved differently after July 2013, we would expect both speculative-rated and unrated firms to behave similarly. The even columns of Table \ref{table_combinedplacebo} show that unrated firms with preferred stock do not behave differently from unrated firms without preferred stock.

Third, we use investment-grade firms with preferred stock as a control group to alleviate concerns that the effect is driven purely by firms with preferred stock, and the results are unchanged (Online Appendix Table \ref{table_diff_c}).

These three robustness checks show that the balance sheet of firms changes after July 2013 only for firms that are affected by the rule change, i.e., firms with preferred stock and rated speculative by Moody's, and not for similar firms with other hybrid securities, or with preferred stock but not rated by Moody's, or with preferred stock but rated investment grade.

Finally, we find similar results again when we use a regular difference-in-differences approach without matching, where the control group includes all speculative-grade firms without preferred stock (Online Appendix Tables \ref{table_debt_all}-\ref{table_bs_all}) and when we use debt over assets as an alternative definition of leverage (Online Appendix Table \ref{table_bla_c}).


\subsection{Market Reactions to Moody's Rule Change} \label{subsec: market_response}

As mentioned in the hypothesis development section, the announcement of the rule change may have revealed new information that speculative-grade firms with preferred stock were less risky than originally thought, or investors simply believed there was new information because of changes in ratings. An alternative explanation is that rating triggers are used in contracts to prevent shareholders from taking actions that destroy firm value when risky debt is in place. We view the increase in investment as consistent with the rating covenant channel; however, we further test these channels by analyzing bond and stock price responses to the rule change.

In Table \ref{table_car}, we include regressions of cumulative abnormal returns (CAR) over various windows around the announcement date of the rule change. To be consistent with the other tests, we always compare the CAR of treated relative to matched control firms. We find statistically significant and positive stock price reactions around the announcement of the rule change, but not before or after the event window. In Columns (1) - (2), we test if there is any run-up in the period of 30 to 2 days prior to the rule change. We find no statistically or economically significant estimates for either a binary or continuous treatment. In Columns (3) - (6), we look at various windows beginning 1 day prior to the announcement up to 1 or 3 days after the announcement. All estimates are positive and statistically significant. For instance, from 1 day prior to 3 days post (Column (5)), we find a 2.8\% positive stock price reaction. This suggests that the equity market viewed the expanded debt capacity as a positive surprise. In Columns (7) - (8), we test if there is a continued run-up after the rule change over the period of 4 days to 30 days after the announcement. The estimates are statistically insignificant for both the binary and continuous treatment. These results also further validate the rule change as being a truly unexpected, exogenous event to firms' prior financing and investment decisions.

In Table \ref{table_spreadsm}, we analyze firms' credit spread responses to the rule change. We obtain daily credit spread data from ICE.\footnote{The benefit of using ICE data over TRACE is that ICE prices do not depend entirely on executed TRACE eligible trades. The prices may reflect quotes from ICE clients, such as investment banks, and trades of bonds that are not TRACE eligible. Furthermore, TRACE often has only a few trades per month, whereas the index constituents' data are updated daily. Although CDS spreads would be ideal, they are not as commonly traded among speculative-grade firms.} In particular, we use constituent bond-level option-adjusted spreads from the ICE BofAML US High Yield Index. To properly compare credit spreads, we exclude all subordinated and secured bonds. For each date, we compute a firm-level average credit spread weighted by each bond's outstanding amount. We are able to find credit spreads for 22 treated firms, as several treated firms only have bank debt, convertible debt, or both, which leads them to be excluded from the index. To increase the power of the test, we perform the same matching on the treated firms but also match on credit spreads on the day before the rule change (July 30, 2013).\footnote{We find similar results with the unmatched regression as shown in Online Appendix Table \ref{table_spreadun}.} We use the same analysis as in Table \ref{table_car}; however, we define our dependent variable as the change in credit spreads (in bps) for each window. In Columns (1) - (2) and (7) - (8) we find no change in spreads before or after the announcement of the rule change, while in Columns (3) - (6), we find an increase in credit spreads that is statistically significant over the window beginning 1 day prior up to 3 days post for both the binary and continuous treatment variable. In Column (5), we find a 16.2bp increase in credit spreads around the same period that stock prices increase by 2.8\%. This compares to the starting average credit spread of 461bps (3.5\% increase).

This negative bond price reaction, coupled with a positive price reaction in the equity market suggests a transfer in value from debtholders to equityholders. Note that as long as the new investment we observe following the rule change is risky, unsecured debtholders will be made worse off. Furthermore, if the new debt is pari-passu or senior to the existing debt, debtholders will also be diluted. In untabulated results, we find no increase in subordinated debt following the rule change; hence, the new debt is at least pari-passu to senior unsecured debt, which would be dilutive to existing unsecured debtholders.\footnote{We obtain subordinated debt from the annual Compustat file, as it is not reported in the quarterly file.}

These asset price responses to the rule change suggest that the increase in debt capacity is caused by the addition of slack in rating-based triggers in financial contracts. If an information-based channel were dominant, we would expect debt prices to increase or remain flat.

It may seem puzzling that firms issue more debt even as their bond yields go up. However, this is a natural consequence of debt markets anticipating future debt issuance that dilutes the value of the firm's existing debt. Indeed, in standard models of dynamic debt issuance, initial yields are higher on debt in anticipation of future debt issuance that dilutes existing debtholders (e.g., \citealp{bizer1992sequential} and \citealp{admati2018leverage}). Despite higher ex-ante yields, firms still have incentives to issue more debt.

\subsection{The Effect of Convertible Debt on Firms' Response to Rule Change}

We next develop additional tests of the covenant channel based on which firms will have the highest incentives to issue debt and take on risky investment following the rule change. \cite{green1984investment} shows that convertible debt reduces the incentives of equityholders to engage in risk-shifting because convertible debtholders benefit from risky investment through their option/equity exposure. Therefore, we examine whether firms with more convertible debt in their capital structure respond less to the rule change. Using the same \textit{Convert} variable defined in Table \ref{table_combinedplacebo}, we test whether firms with more convertible debt respond less to the rule change.\footnote{The convertible debt field is only available in the Compustat annual file; hence, we take the last ratio of convertible debt to total debt prior to the rule change. We find similar results when using convertible debt over assets rather than over total debt.} 21\% of firms in our sample have convertible debt that accounts for 22\% of these firms' total debt. Furthermore, the use of convertible debt is similar between treated firms and control firms: 27\% of treated firms have convertible debt, compared to 25\% of matched control firms and 20\% of all other speculative-rated firms. The results are displayed in Table \ref{table_convert}, which shows that treated firms issue less debt and invest less if they have convertible debt at the time of the rule change. For example, in Column (1), a one percentage point increase in convertible debt results in treated firms increasing their debt by 1.5\% less following the rule change. We find similar results when we use as counterfactual investment-grade firms with preferred stock (see Table \ref{table_convertIG}). We also re-estimate Table \ref{table_convert} using a dummy variable equal to one if a firm had any positive convertible debt at the time of the rule change. The results, reported in Table \ref{table_convert_dummy}, are virtually unchanged in both economic magnitude and statistical significance. This additional test reinforces that our findings regarding the role of convertible debt are not driven by outliers.

Overall, these results are consistent with convertible debt reducing firms' incentives to increase their risk at the expense of debtholders and are difficult to explain with an information-based channel of rating influence. We argue that our results are consistent with the hypothesis that the rule change relaxed covenants, as firms increase their investment as their liabilities become riskier, and this response is stronger the greater their incentives to risk-shift. \cite{gilje2016firms} finds evidence that financial covenants are used ex-ante to prevent risk-shifting. Our results are consistent with \cite{gilje2016firms} in that once rating-based triggers are relaxed, firms with risk-shifting incentives increase their risky investment at the expense of existing debtholders.

However, we caution that the sample of treated firms with convertible debt is small (only 12 firms), and therefore, we interpret these results as suggestive evidence consistent with a risk-shifting channel rather than definitive proof that this channel dominates. This test is one piece of a broader set of evidence, which includes stock price and bond spread responses, that collectively point toward covenants as the operative mechanism.


To see if firms are indeed risk-shifting following the rule change, we could test whether total firm value decreases. Unfortunately, we only have market prices and maturities for a small subset of the firms' liabilities. To truly test if firm value decreases, we would need to know the change in value of all of the firms' liabilities, including derivatives, leases, pensions, junior debt, etc.\footnote{In particular, the treated firms' preferred stock likely changed value substantially given their perpetual nature.}  Although we have strong evidence of a transfer from debtholders to equityholders, we cannot definitively claim that the overall value of the firms decreased. In addition, given that we have only one event, it is possible that ex-ante covenants were put in place to prevent risk-shifting, but ex-post, they are too tight and relaxing them increases firm value.


We could also test the contracting channel by looking at the presence of rating triggers directly. Unfortunately, there is a vast underreporting of rating triggers. While \cite{moodytrigger2} found that 87.5\% of US corporate issuers have ratings triggers with an average of over four triggers per company, only 22.5\% of the triggers were disclosed in SEC filings\footnote{In the Online Appendix, we include excerpts from treated firms' filings that explicitly reference contractual ties to credit ratings.}. Furthermore, firms may have an incentive to hide rating triggers exactly when they are most important, which could bias an analysis that focuses on disclosed triggers.\footnote{For instance, Enron had many rating triggers that market participants were not aware of that hastened its demise. See \href{https://www.nytimes.com/2002/01/22/business/enron-s-collapse-credit-ratings-enron-spurs-debt-agencies-consider-some-changes.html}{Enron's Collapse: Credit Ratings; Enron Spurs Debt Agencies To Consider Some Changes}} Consistent with the observed underreporting, using DealScan and FISD, we find three treated firms that disclosed performance-based rating triggers in their loan contracts and zero with bond rating triggers. There also exist rating triggers in bank loans that induce puts or force firms to post more collateral from rating downgrades, which do not appear in DealScan (\citealp{moodytrigger2}). In addition, there may be many other forms of triggers outside of bank and bond triggers. For instance, a large number of triggers are found in lease agreements, commercial agreements and derivative contracts (\citealp{moodytrigger2}). Even knowing that a trigger exists may not be sufficient for our analysis because we do not know the terms of the trigger (e.g., below what rating firms must post more collateral, and how much collateral they must post). We believe our work further highlights the importance of disclosing the existence and terms of rating triggers, as they appear to have a material effect on firms' behavior.


\subsection{Moody's Rule Change and Firms' Security Design}\label{sec:secdesign}

We have shown that Moody's rule change caused treated firms to increase their debt levels, leverage, and investment. However, the rule change had a broader impact beyond firms with existing preferred stock in their capital structure. The new methodology treats preferred stock more favorably for all speculative-grade firms. Thus, we should expect firms rated speculative-grade by Moody's to be more likely to increase their preferred stock holdings on their balance sheets after the rule change. Unfortunately, the new methodology applies to all firms rated speculative by Moody's, thereby limiting the options for choosing a natural counterfactual. We thus consider all firms rated speculative-grade only by Moody's as treated, and, as a control group, firms rated speculative-grade by both S\&P and Moody's, or just by S\&P. At the time of the rule change, there were 66 firms rated speculative-grade only by Moody's. The firms only rated speculative by Moody's represent only a subset of the firms that should be affected by the rule change, but it helps with identification, as one group clearly has a higher incentive to issue preferred stock than the other, since S\&P did not change its preferred stock methodology. We expect that Moody's-only rated speculative-grade firms should have more preferred stock after the rule change compared to other speculative-grade firms.\footnote{We find similar results when we compare firms only rated speculative-grade by Moody's to firms only rated speculative-grade by S\&P, but have less statistical power.}

Table \ref{table_issue} shows the coefficients of a generalized difference-in-differences where we compare the percent of preferred stock relative to capital and simply the total amount of preferred between Moody's-only-junk firms and all other speculative-grade firms, before versus after the Moody's rule change. The sample is restricted to speculative-grade firms in Columns (1) and (2). We find that the use of preferred stock increases after the rule change for Moody's-only junk firms. The coefficient is not only statistically significant, but it is also economically large: Before the rule change, preferred stock accounted for only 0.33\% of total capital for speculative-grade firms rated only by Moody's, and the amount of preferred stock relative to capital almost triples after the rule change. As a placebo test, we perform the same analysis on firms rated investment-grade. We would expect no change in the use of preferred stock, given that the new rule did not apply to investment-grade firms. In Columns (3) and (4), we indeed find no change in the use of preferred stock by firms only rated investment-grade by Moody's.

This analysis shows the rule change has a pervasive effect, not only on firms with existing preferred stock, but on all firms rated speculative by Moody's. Thus, the rule change has two main effects: (i) a direct one that exogenously lowered the leverage of speculative firms with preferred stock that were rated by Moody's; and (ii) an indirect effect that increased the advantage of issuing preferred stock to all speculative firms rated by Moody's. One might then worry that firms with no preferred stock at the announcement are not a valid control group in the main analysis because some of them (those rated speculative-grade by Moody's) are also indirectly affected by the rule change. However, the effect of the rule change on new issuance of preferred stock applies equally to both treated and control groups. Also, as shown in Table \ref{table_sum}, the proportion of firms rated speculative-grade only by Moody's is statistically indistinguishable across our main treated and control groups.

Alternatively, one could argue that the proper treatment group for our entire analysis is firms rated speculative-grade only by Moody's. However, the current empirical approach allows us to measure the direct and indirect effects separately, thus showing heterogeneity in the effect across speculative-rated firms. In fact, the direct effect appears to be much larger than the indirect effect. Therefore, we view the dominant effect as coming from firms with existing preferred stock in their capital structure.


One may also wonder why firms do not simply all issue preferred stock rather than more debt. First, rating covenants are often asymmetric, so the benefit of improving the rating may not be as much as the cost of being downgraded.\footnote{Many of these rating triggers occur only when the firm's rating drops below a specific level.} This idea is also consistent with firms targeting a \textit{minimum} rating (e.g., \citealp{hovakimian2009credit}). Second, preferred stock is not a very liquid security, and hence, even with this improved equity credit, there are likely limits to how much preferred stock firms would issue.

Our evidence of increased preferred stock issuance is consistent with anecdotal evidence suggesting the rule change affected the design of hybrid securities (\citealp{wsj}). These results also relate to \cite{rauh2010capital}, who show that there is substantial heterogeneity in the types of debt firms use. Some of the observed heterogeneity in debt structure may be due to rating agency methodologies, as firms adjust their security design decisions based on how those securities are treated by rating agencies. To our knowledge, this is the first paper to show that rating agency methodologies have a causal effect on firms' security design decisions.


\subsection{Moody's Rule Change and Credit Ratings}

The change in the treatment of preferred stock mechanically lowered the leverage, as assessed by Moody's, of the affected firms. This reduction in leverage amounts to approximately a one-notch improvement in the credit ratings. However, firms responded to such an increase in debt capacity by issuing more debt. We would thus not expect any change in firms' ratings. Consistent with this hypothesis, Columns (1) and (2) of Table \ref{table_cr} show that Moody's ratings did not significantly change after the new methodology, consistent with firms considering leverage ratios as defined by rating agencies.

However, 89\% of the treated firms are rated both by S\&P and Moody's, and S\&P did not alter its methodology. If firms issued more debt in response to the additional debt capacity allowed by Moody's, then we should observe a higher likelihood of S\&P downgrading the treated firms. Columns (3) and (4) of Table \ref{table_cr} show that the S\&P ratings do not change significantly between treated and control firms after the rule change. A possible explanation of this finding is that the response of firms to Moody's rule change is heterogeneous: firms that are also rated by S\&P respond less to the rule change because they are concerned that increasing debt might lead to a downgrade by S\&P. We test this hypothesis using a triple-difference specification, and thus comparing the effect of the rule change for firms only rated by Moody's, and firms that are rated both by Moody's and S\&P. Table \ref{table_cross} shows that firms rated by S\&P appear to respond less, although the statistical significance of a triple-difference test is weak given the small number of firms only rated by Moody's. This result suggests that firms have a joint problem in targeting a leverage ratio that may depend on multiple rating agencies. Having multiple ratings also seems to mitigate the effect of rating agency-specific methodologies on firms' financing and investment decisions.

If firms rated by S\&P respond less, we would expect their Moody's ratings to improve following the rule change. In Column (4) of Table \ref{table_cross}, we estimate the same regressions except that we include Moody's rating as the dependent variable. We find some evidence supporting this, as the coefficient on the triple-difference estimate is negative, though it is not statistically significant. Although we do not have sufficient statistical power to show that this is definitively occurring, the coefficient's direction is consistent with a heterogeneous response among treated firms depending on whether they are rated by S\&P.

Although most firms are rated by both S\&P and Moody's, many firms also rated by S\&P could have had room to increase their leverage without being downgraded, which explains why we see a response even among firms rated by both rating agencies.

One could argue that because we see no change in Moody's ratings, the rule change had no real effect on Moody's risk assessments. However, we find ample evidence that suggests otherwise. First, affected firms dramatically changed their balance sheets after the rule change. Second, there were economically meaningful asset price responses around the announcement of the rule change. Third, firms rated only speculative-grade by Moody's increased the amount of preferred stock in their capital structure to take advantage of the new methodology. Furthermore, the fact that ratings did not change is consistent with firms targeting a minimum credit rating.


\section{Conclusion}

This paper shows that rating agency methodologies have a causal effect on firms' financing, investment, and security design decisions. We also provide evidence that ratings are an important ex-ante contracting tool to mitigate shareholder and debtholder conflicts.

A methodological change by Moody's to the equity credit assigned to preferred stock led to a large and persistent increase in debt levels and leverage for speculative-grade firms with preferred stock in their capital structure. Our evidence suggests that the rule change created debt capacity, and firms use the proceeds from the new debt to increase their asset base through investment, rather than simply recapitalizing. Stock prices of treated firms increased while bond prices decreased, suggesting that the rule change increased shareholder value at the detriment of existing debtholders. Furthermore, firms with convertible debt, which theory predicts are less subject to risk-shifting incentives, respond less to the rule change. Finally, rating agencies causally influence firms' security design problem: we find that preferred stock usage nearly tripled after the methodology change. Collectively, our evidence is consistent with the rule change relaxing financial constraints through rating-based covenants. We view our results as the first experiment in the literature that provides tests to distinguish between information and contract-based channels of rating agency influence.


The findings of the paper inform firms, investors, rating agencies, and regulators about the unintended consequences of even seemingly small changes in methodologies used by rating agencies. It also further highlights the importance of ratings and rating agency actions for firms and financial markets.

\doublespacing


\clearpage
\newpage

\bibliography{bib.bib}
	\bibliographystyle{jfnew}

\clearpage


\newpage

\begin{figure}[bp!]
\begin{center}
\includegraphics[width=.9\textwidth]{"Prefeq_Histogram".png}
\vspace{2mm}
\end{center}
\begin{singlespace}
\begin{small}
\caption{\textbf{Preferred Stock Histogram}}
\label{fig_histo}
\noindent Figure \ref{fig_histo} includes a histogram of Preferred/Capital, the ratio between the amount of preferred stock and the sum of debt and shareholders' equity in the last quarter prior to July 2013, when Moody's changed its debt adjustment methodology. Leverage, as calculated by Moody's, immediately dropped by half of this number at the time of the rule change.
\end{small}
\end{singlespace}
\end{figure}
\clearpage
\newpage

\begin{figure}[bp!]
Total Debt
\begin{center}
\includegraphics[width=1\textwidth]{"Logdebt TS Continous Matched_medium".png}
\end{center}


\begin{singlespace}
\begin{small}
\caption{\textbf{Diff-in-Diff Time Trends of Debt Levels}}
\label{fig_debt_reg}
\noindent Figure \ref{fig_debt_reg} plots regression coefficients of $Preferred/Capital \times \lambda_{t,k}$ with 90\% confidence intervals estimated from Equation (\ref{coefficients}). Controls include: profitability, tangibility, sales and market-to-book. The sample period is 2012Q2 - 2015Q2 and includes treated and matched control firms. The dependent variable is Total Debt. Standard errors are clustered by firm.

\end{small}
\end{singlespace}
\end{figure}
\clearpage
\newpage

\begin{figure}[bp!]


Lev\textsubscript{Moody's}
\begin{center}
\includegraphics[width=1\textwidth]{"RL TS Continous Matched_medium".png}
\end{center}
\begin{singlespace}
\begin{small}
\caption{\textbf{Diff-in-Diff Time Trends of Leverage Ratios}}
\label{fig_bl_reg}
\noindent Figure \ref{fig_bl_reg} plots regression coefficients of $Preferred/Capital \times \lambda_{t,k}$ with 90\% confidence intervals estimated from Equation (\ref{coefficients}). Controls include: profitability, tangibility, sales and market-to-book. The sample period is 2012Q2 - 2015Q2 and includes treated and matched firms. The dependent variable is Lev\textsubscript{Moody's}. Standard errors are clustered by firm.
\end{small}
\end{singlespace}
\end{figure}
\clearpage
\newpage


\begin{figure}[bp!]
Panel A: PP\&E
\begin{center}
\includegraphics[width=1\textwidth]{"PPE TS Continous Matched_medium".png}
\vspace{2mm}
\end{center}
Panel B: Assets
\begin{center}
\includegraphics[width=1\textwidth]{"Assets TS Continous Matched_medium".png}
\end{center}
\begin{singlespace}
\begin{small}
\caption{\textbf{Diff-in-Diff Time-Trends of Assets Levels}}
\label{fig_bs_reg}
\noindent Figure \ref{fig_bs_reg} plots regression coefficients of $Preferred/Capital \times \lambda_{t,k}$ with 90\% confidence intervals estimated from Equation (\ref{coefficients}). Controls include: profitability, tangibility, sales and market-to-book. The sample period is 2012Q2 - 2015Q2 and includes treated and matched firms. In Panel A the dependent variable is PP\&E and in Panel B the dependent variable is Assets. Standard errors are clustered by firm.
\end{small}
\end{singlespace}
\end{figure}
\clearpage
\newpage


\begin{table}[htbp]
\caption{\textbf{Variable Definitions}
} \label{table_vardef}

\begin{singlespace}
\begin{small}
This table identifies the data sources and describes the construction of variables used in the analysis. Quarterly financial data is from Compustat, ratings data is from Thomson Eikon, stock returns are from CRSP and credit spreads are from ICE.
\end{small}
\end{singlespace}
\vspace{2mm}

\begin{footnotesize}
\renewcommand{\arraystretch}{1.5}
\begin{tabular}{@{}lll@{}}

\hline
\textbf{Variable} & \textbf{Definition} & \textbf{Comments} \\
\hline
Total Debt & log total debt [log(dlcq + dlttq)] &  \\
Long-Term Debt & log long-term debt [log(dlttq)] &  \\
Short-Term Debt & log short-term debt [log(dlcq)]&  \\
Lev\textsubscript{GAAP} & total debt [dlcq + dlttq] \\
& / (total debt + equity [seqq (teqq if missing seqq)]) & Winsorized [0,1]\\
Lev\textsubscript{Moody's} &

\begin{math}
=
\begin{cases}
 \frac{\textnormal{total debt +0.5 preferred}}{\textnormal{total debt + equity}} & \textnormal{ before July 31st, 2013 } \\
\frac{\textnormal{total debt}}{\textnormal{total debt + equity}}  & \textnormal{ after July 31st, 2013 }
\end{cases}
\end{math}

& Winsorized [0,1]\\
Debt/Assets\textsubscript{GAAP} & total debt [dlcq + dlttq] / total assets [atq] & Winsorized [0,1] \\
Debt/Assets\textsubscript{Moody's} &

\begin{math}
=
\begin{cases}
 \frac{\textnormal{total debt + 0.5 preferred}}{\textnormal{assets}} & \textnormal{ before July 31st, 2013 } \\
 \frac{\textnormal{total debt }}{\textnormal{assets }}   & \textnormal{ after July 31st, 2013 }
\end{cases}
\end{math}
& Winsorized [0,1]\\

Preferred &   preferred\textsubscript{t=0} / 1000 & \\
Log(Preferred) &   log(1+preferred)  & \\
Preferred/Capital &   preferred\textsubscript{t=0} / (total debt + equity) & Winsorized [0,1]\\
Market Equity & total shares outstanding [cshoq/1000] $\times$  price [prccq] \\


Assets & log assets [log(atq)] &  \\
PPE & log ppe [log(ppentq)] & \\
Capex & log capex [log(1+capxq)]&  \\
Market-to-Book & (market equity + total debt + preferred [pstkq] + \\
& \quad deferred taxes [txditcq]) / total assets [atq] & Winsorized [1,99] \\
Tangibility & property plant \& equipment [ppentq] / total assets [atq]  & Winsorized [1,99] \\
Profitability & EBITDA [oibdpq] / total assets [atq]  & Winsorized [1,99] \\
Sales & log sales [log(1+saleq)]  & Winsorized [1,99] \\
Convert & convertible debt\textsubscript{t=0} (dcvt) / total debt\textsubscript{t=0} (dlc + dltt)  & Winsorized [0,1], from annual file \\
Moody's rating & 1 - 22 corresponding to letter ratings (highest to lowest)  &  \\
S\&P rating & 1 - 22 corresponding to letter ratings (highest to lowest)  &   \\


SP &  \begin{math}
=
\begin{cases}
 1  & \textnormal{ if rated by S\&P on July 31st 2013} \\
0 & \textnormal{ otherwise }
\end{cases}
\end{math}  \\
MJOnly &\begin{math}
=
\begin{cases}
 1  & \textnormal{ if only rated junk by Moody's on July 31st 2013 } \\
0 & \textnormal{ otherwise}

\end{cases}
\end{math}  \\
MIGOnly & \begin{math}
=
\begin{cases}
 1  & \textnormal{ if only rated IG by Moody's on July 31st 2013 } \\
0 & \textnormal{  otherwise}
\end{cases}
\end{math} \\
CAR[-1,1] & sum of the firm's daily stock return minus the \\  &  value-weighted CRSP market return from $t-1$ to $t+1$   &   \\
Credit Spread & firm's average option-adjusted spread, weighted by &   \\ &
amount outstanding  \\
$\Delta CS[-1,1]$ & change in credit spread from $t-1$ to $t+1$ \\
\hline

\end{tabular}
\end{footnotesize}
\end{table}
\clearpage
\newpage


\begin{landscape}
\pagestyle{empty}
\begin{table}[htbp]
\caption{\textbf{Summary Statistics}
} \label{table_sum}
\begin{singlespace}
\begin{small}

\noindent This table compares treated firms, all untreated speculative-grade firms and matched untreated speculative-grade firms at the time of Moody's rule change. All variable definitions are included in the Table \ref{table_vardef}. Long-term debt, short-term debt, PPE, sales and assets are displayed in billions of dollars in this table (rather than logs as in the regressions). *, **, and *** indicate statistical significance at the $10\%$, $5\%$, and $1\%$ levels, respectively, based on a t-test.

\end{small}
\end{singlespace}
\vspace{3mm}
\small
  \centering
{
\def\sym#1{\ifmmode^{#1}\else\(^{#1}\)\fi}
\begin{tabular}{l*{3}{ccccccccccc}}
\hline
                    &\multicolumn{4}{c}{Treated Firms}               &\multicolumn{5}{c}{All Untreated Speculative-Grade Firms}
                    &\multicolumn{5}{c}{Matched Speculative-Grade Firms}   \\
            & & & & & & & & & Diff w.r.t & & & & & Diff w.r.t
            \\
                      &           N&        Mean&      Median&          SD&           N&        Mean&      Median&          SD& Treated         &           N&        Mean&      Median&          SD&  Treated         \\
\hline
Preferred Stock     &          44&       0.395&        0.10&       1.112&         431&       0.000&        0.00&       0.000&      -0.395\sym{**} &         164&       0.000&        0.00&       0.000&      -0.395\sym{**} \\
Preferred/Capital   &          44&       0.096&        0.07&       0.094&         431&       0.000&        0.00&       0.000&      -0.096\sym{***}&         164&       0.000&        0.00&       0.000&      -0.096\sym{***}\\
Book Leverage       &          44&       0.571&        0.54&       0.232&         431&       0.588&        0.55&       0.248&       0.017         &         164&       0.571&        0.53&       0.210&       0.000         \\
Long-Term Debt      &          44&       2.422&        0.97&       3.895&         431&       1.929&        0.85&       3.496&      -0.493         &         164&       2.497&        0.87&       5.721&       0.075         \\
Short-Term Debt     &          44&       0.323&        0.01&       1.501&         431&       0.146&        0.01&       0.586&      -0.178         &         164&       0.185&        0.00&       0.647&      -0.138         \\
Market Equity       &          44&       3.053&        1.04&       7.170&         417&       2.988&        1.84&       4.239&      -0.066         &         164&       2.894&        1.14&       6.407&      -0.159         \\
Market to Book      &          44&       0.930&        0.91&       0.311&         417&       1.241&        1.12&       0.537&       0.311\sym{***}&         164&       0.988&        0.92&       0.266&       0.058         \\
Moody's Rating       &          44&      15.227&       15.00&       2.078&         431&      14.188&       14.00&       1.914&      -1.039\sym{***}&         164&      14.817&       15.00&       1.749&      -0.410         \\
S\&P Rated          &          44&       0.886&        1.00&       0.321&         431&       0.919&        1.00&       0.273&       0.032         &         164&       0.945&        1.00&       0.228&       0.059         \\
S\&P Rating         &          39&      14.333&       15.00&       1.782&         396&      13.240&       13.00&       1.796&      -1.093\sym{***}&         155&      13.794&       14.00&       1.557&      -0.540\sym{*}  \\
PPE                 &          44&       3.384&        0.80&       7.666&         422&       2.006&        0.76&       3.870&      -1.378         &         164&       2.951&        1.11&       6.274&      -0.433         \\
Profitability       &          44&       0.012&        0.01&       0.014&         425&       0.019&        0.02&       0.021&       0.007\sym{***}&         164&       0.013&        0.02&       0.014&       0.001         \\
Tangibility         &          44&       0.445&        0.44&       0.308&         422&       0.377&        0.33&       0.276&      -0.068         &         164&       0.454&        0.42&       0.307&       0.010         \\
Sales               &          44&       1.831&        0.39&       5.932&         431&       0.944&        0.44&       1.453&      -0.887         &         164&       0.986&        0.31&       1.813&      -0.845         \\
Assets              &          44&       9.645&        2.51&      25.653&         425&       4.881&        2.48&       7.833&      -4.764         &         164&       6.076&        2.48&      12.659&      -3.569         \\
\hline
Observations        &          44&            &            &            &         431&            &            &            &         475         &         164&            &            &            &         208         \\
\hline
\end{tabular}
}

\end{table}
\end{landscape}


\clearpage

\newpage

\begin{table}[htbp]
\caption{\textbf{The Effect of Rule Change on Debt Levels}} \label{table_debt}
\begin{singlespace}
\begin{small}

\noindent This table contains results testing whether treated firms increase their debt levels after the rule change compared to matched firms. Coefficients are estimated using Equation (\ref{eq:single}). The dependent variables are total debt (Columns 1 and 2), long-term debt (Columns 3 and 4) and short-term debt (Columns 5 and 6) and are all in logs. Preferred Dummy is an indicator variable that equals one if the firm had preferred stock in its capital structure in the last quarter prior to the rule change. Preferred/Capital is the ratio between the amount of preferred stock and the sum of debt and shareholders' equity in the last quarter prior to the rule change. T-statistics are shown below the parameter estimates in parentheses and are calculated using robust standard errors clustered by firm. *, **, and *** indicate statistical significance at the $10\%$, $5\%$, and $1\%$ levels, respectively.

\end{small}
\end{singlespace}
\vspace{3mm}
\small
  \centering
    \begin{tabular}{lcccccc}
    \toprule
    \textbf{} &
      \multicolumn{2}{c}{\textbf{Total Debt}} &  \multicolumn{2}{c}{\textbf{LT Debt}} &
      \multicolumn{2}{c}{\textbf{ST Debt}}
      \\
      \cline{2-7}
            &\multicolumn{1}{c}{(1)}   &\multicolumn{1}{c}{(2)}   &\multicolumn{1}{c}{(3)}   &\multicolumn{1}{c}{(4)}   &\multicolumn{1}{c}{(5)}   &\multicolumn{1}{c}{(6)}   \\
\hline
Preferred Dummy x Post&       0.223***&               &       0.242***&               &       0.177   &               \\
            &      (3.53)   &               &      (3.64)   &               &      (0.61)   &               \\
Preferred/Capital x Post&               &       1.267***&               &       1.440***&               &       2.859   \\
            &               &      (3.30)   &               &      (3.19)   &               &      (1.10)   \\
Profitability&      -0.008   &      -0.021   &      -0.749** &      -0.794** &      -1.550   &      -1.515   \\
            &     (-0.03)   &     (-0.09)   &     (-2.03)   &     (-2.08)   &     (-1.46)   &     (-1.41)   \\
Tangibility &       0.209   &       0.210   &      -0.028   &      -0.032   &      -0.219   &      -0.208   \\
            &      (0.46)   &      (0.44)   &     (-0.06)   &     (-0.06)   &     (-0.13)   &     (-0.12)   \\
Sales       &       0.357***&       0.364***&       0.440***&       0.447***&      -0.202   &      -0.222   \\
            &      (5.82)   &      (5.98)   &      (5.57)   &      (5.68)   &     (-0.89)   &     (-0.98)   \\
Market-to-Book&      -0.036   &      -0.041   &      -0.063   &      -0.073   &      -0.301   &      -0.316   \\
            &     (-0.19)   &     (-0.21)   &     (-0.35)   &     (-0.39)   &     (-0.80)   &     (-0.84)   \\
\_cons      &       4.826***&       4.806***&       4.412***&       4.401***&       4.933***&       5.060***\\
            &     (10.46)   &     (10.30)   &      (7.19)   &      (7.15)   &      (3.31)   &      (3.37)   \\
            &               &               &               &               &               &               \\
Firm FE     &           Y   &           Y   &           Y   &           Y   &           Y   &           Y   \\
Quarter x Cohort FE&           Y   &           Y   &           Y   &           Y   &           Y   &           Y   \\
Firm Quarters&        2448   &        2448   &        2425   &        2425   &        1730   &        1730   \\
R\textsuperscript{2}&       0.128   &       0.114   &       0.149   &       0.137   &       0.007   &       0.009   \\



    \bottomrule
    \end{tabular}
\end{table}
\clearpage

\newpage


\begin{landscape}
\pagestyle{empty}
\begin{table}[htbp]
\caption{\textbf{Yearly Placebo Tests}} \label{table_placebo}
\begin{singlespace}
\begin{small}

This table contains results testing whether there are similar changes in total debt in the second half of other years in which the rule change did not take place among treated and matched firms. Coefficients are estimated using Equation (\ref{eq:single}).  Preferred Dummy is an indicator variable that equals one if the firm had preferred stock in its capital structure in the last quarter prior to the rule change. After July is a dummy variable that equals 1 if the quarter end is after July. Controls include: profitability, tangibility, sales and market-to-book. T-statistics are shown below the parameter estimates in parentheses and are calculated using robust standard errors clustered by firm. *, **, and *** indicate statistical significance at the $10\%$, $5\%$, and $1\%$ levels, respectively
\end{small}
\end{singlespace}
\vspace{3mm}

\small
  \centering
\begin{tabular}{lcccccccc}
    \toprule
    \textbf{} &
      \multicolumn{8}{c}{\textbf{Total Debt}}
      \\
                \cline{2-9}
              &\multicolumn{1}{c}{(1)}   &\multicolumn{1}{c}{(2)}   &\multicolumn{1}{c}{(3)}   &\multicolumn{1}{c}{(4)}   &\multicolumn{1}{c}{(5)}   &\multicolumn{1}{c}{(6)}   &\multicolumn{1}{c}{(7)}   &\multicolumn{1}{c}{(8)}   \\
\hline
Preferred x After July&       0.037   &       0.049   &      -0.138   &      -0.025   &      -0.029   &       0.132***&       0.043   &       0.081   \\
            &      (0.48)   &      (0.85)   &     (-1.48)   &     (-0.60)   &     (-0.69)   &      (3.03)   &      (1.04)   &      (1.15)   \\
Constant    &       5.591***&       6.680***&       3.680*  &       5.665***&       5.501***&       6.847***&       6.730***&       4.757***\\
            &     (12.73)   &      (8.15)   &      (1.97)   &      (7.79)   &      (8.39)   &     (11.62)   &     (17.18)   &      (3.22)   \\
            &               &               &               &               &               &               &               &               \\
Year        &        2008   &        2009   &        2010   &        2011   &        2012   &        2013   &        2014   &        2015   \\
Controls    &           Y   &           Y   &           Y   &           Y   &           Y   &           Y   &           Y   &           Y   \\
Firm FE     &           Y   &           Y   &           Y   &           Y   &           Y   &           Y   &           Y   &           Y   \\
Quarter x Cohort FE&           Y   &           Y   &           Y   &           Y   &           Y   &           Y   &           Y   &           Y   \\
Firm Quarters&         695   &         707   &         742   &         800   &         825   &         823   &         814   &         760   \\
R\textsuperscript{2}&       0.068   &       0.006   &       0.035   &       0.392   &       0.452   &       0.222   &       0.050   &       0.053   \\


    \bottomrule
    \end{tabular}
\end{table}
\end{landscape}
\clearpage
\newpage


\begin{table}[htbp]
\caption{\textbf{The Effect of Rule Change on Leverage}
} \label{table_lev}
\begin{singlespace}
\begin{small}

\noindent This table contains results testing whether treated firms increase their GAAP leverage after the rule change compared to matched firms. Coefficients are estimated using Equation (\ref{eq:single}). The dependent variables are Lev\textsubscript{GAAP} and Levtextsubscript{Moody's}. Preferred Dummy is an indicator variable that equals one if the firm had preferred stock in its capital structure in the last quarter prior to the rule change. Preferred/Capital is the ratio between the amount of preferred stock and the sum of debt and shareholders' equity in the last quarter prior to the rule change. T-statistics are shown below the parameter estimates in parentheses and are calculated using robust standard errors clustered by firm. *, **, and *** indicate statistical significance at the $10\%$, $5\%$, and $1\%$ levels, respectively.

\end{small}
\end{singlespace}
\vspace{3mm}
\small
  \centering
   \scalebox{1}{ \begin{tabular}{lcccc}
    \toprule
    \textbf{ } &
      \multicolumn{2}{c}{\textbf{Lev\textsubscript{GAAP}}}  &
        \multicolumn{2}{c}{\textbf{Lev\textsubscript{Moody's}}}
      \\
          \cline{2-5}

              &\multicolumn{1}{c}{(1)}   &\multicolumn{1}{c}{(2)}   &\multicolumn{1}{c}{(3)}   &\multicolumn{1}{c}{(4)}   \\
\hline
Preferred Dummy x Post&       0.031*  &               &      -0.016   &               \\
            &      (1.69)   &               &     (-0.86)   &               \\
Preferred/Capital x Post&               &       0.225   &               &      -0.313** \\
            &               &      (1.60)   &               &     (-2.05)   \\
Profitability&      -0.217** &      -0.220** &      -0.236** &      -0.229** \\
            &     (-1.99)   &     (-2.02)   &     (-2.10)   &     (-2.07)   \\
Tangibility &       0.367** &       0.369** &       0.369** &       0.361** \\
            &      (2.38)   &      (2.38)   &      (2.41)   &      (2.34)   \\
Sales       &      -0.011   &      -0.011   &      -0.014   &      -0.012   \\
            &     (-0.53)   &     (-0.51)   &     (-0.66)   &     (-0.59)   \\
Market-to-Book&      -0.070** &      -0.072** &      -0.079** &      -0.074** \\
            &     (-1.99)   &     (-2.00)   &     (-2.17)   &     (-2.06)   \\
\_cons      &       0.553***&       0.552***&       0.585***&       0.578***\\
            &      (4.35)   &      (4.38)   &      (4.64)   &      (4.57)   \\
            &               &               &               &               \\
Firm FE     &           Y   &           Y   &           Y   &           Y   \\
Quarter x Cohort FE&           Y   &           Y   &           Y   &           Y   \\
Firm Quarters&        2454   &        2454   &        2454   &        2454   \\
R\textsuperscript{2}&       0.050   &       0.049   &       0.056   &       0.065   \\



    \bottomrule
    \end{tabular}}
\end{table}


\clearpage

\newpage


\begin{table}[htbp]
\caption{\textbf{The Effect of Rule Change on the Balance Sheet}} \label{table_bs}
\begin{singlespace}
\begin{small}

\noindent This table contains results testing whether treated firms increase the size of their balance sheets after the rule change compared to matched firms. Coefficients are estimated using Equation (\ref{eq:single}). The dependent variables are Capex, PP\&E, and Assets, all in logs. Preferred Dummy is an indicator variable that equals one if the firm had preferred stock in its capital structure in the last quarter prior to the rule change. Preferred/Capital is the ratio between the amount of preferred stock and the sum of debt and shareholders' equity in the last quarter prior to the rule change. T-statistics are shown below the parameter estimates in parentheses and are calculated using robust standard errors clustered by firm. *, **, and *** indicate statistical significance at the $10\%$, $5\%$, and $1\%$ levels, respectively.

\end{small}
\end{singlespace}
\vspace{3mm}
\small
  \centering
    \begin{tabular}{lcccccc}
    \toprule
    \textbf{ } &
      \multicolumn{2}{c}{\textbf{Capex}} &  \multicolumn{2}{c}{\textbf{PPE}} &
      \multicolumn{2}{c}{\textbf{Assets}}
      \\
      \cline{2-7}
              &\multicolumn{1}{c}{(1)}   &\multicolumn{1}{c}{(2)}   &\multicolumn{1}{c}{(3)}   &\multicolumn{1}{c}{(4)}   &\multicolumn{1}{c}{(5)}   &\multicolumn{1}{c}{(6)}   \\
\hline
Preferred Dummy x Post&       0.103   &               &       0.081** &               &       0.083** &               \\
            &      (1.38)   &               &      (2.31)   &               &      (2.23)   &               \\
Preferred/Capital x Post&               &       0.874*  &               &       0.473*  &               &       0.563*  \\
            &               &      (1.83)   &               &      (1.74)   &               &      (1.81)   \\
Profitability&       0.017   &       0.002   &       0.332** &       0.327** &       0.449***&       0.442***\\
            &      (0.05)   &      (0.01)   &      (2.47)   &      (2.39)   &      (3.56)   &      (3.43)   \\
Tangibility &       1.678** &       1.691** &       1.966***&       1.967***&      -0.410   &      -0.407   \\
            &      (2.37)   &      (2.38)   &      (4.67)   &      (4.63)   &     (-1.24)   &     (-1.22)   \\
Sales       &       0.402***&       0.403***&       0.352***&       0.354***&       0.339***&       0.341***\\
            &      (4.69)   &      (4.70)   &      (7.38)   &      (7.48)   &      (6.71)   &      (6.78)   \\
Market-to-Book&       0.303** &       0.296** &      -0.144** &      -0.146** &      -0.139** &      -0.143** \\
            &      (2.39)   &      (2.26)   &     (-2.21)   &     (-2.22)   &     (-2.07)   &     (-2.11)   \\
\_cons      &       0.123   &       0.123   &       3.886***&       3.879***&       6.275***&       6.271***\\
            &      (0.21)   &      (0.21)   &     (13.39)   &     (13.29)   &     (18.98)   &     (18.88)   \\
            &               &               &               &               &               &               \\
Firm FE     &           Y   &           Y   &           Y   &           Y   &           Y   &           Y   \\
Quarter x Cohort FE&           Y   &           Y   &           Y   &           Y   &           Y   &           Y   \\
Firm Quarters&        2450   &        2450   &        2454   &        2454   &        2454   &        2454   \\
R\textsuperscript{2}&       0.075   &       0.075   &       0.337   &       0.334   &       0.302   &       0.299   \\



    \bottomrule
    \end{tabular}
\end{table}
\clearpage

\newpage


\begin{landscape}
\pagestyle{empty}
    \begin{table}[htbp]
\caption{\textbf{Placebo Tests: Alternate Treatment Groups}} \label{table_combinedplacebo}
\begin{singlespace}
\begin{small}

\noindent This table contains results of placebo tests on the effect of the rule change on firms' balance sheets. In columns (1), (3), (5), and (7), we consider the same sample of speculative-rated firms, but use convertible debt instead of preferred stock as treatment variable. In columns (2), (4), (6), (8), we use the sample of unrated firms, instead of firms rated speculative-grade by Moody's. Convert is the ratio of convertible debt to total debt at the time of the rule change.  T-statistics are shown below the parameter estimates in parentheses and are calculated using robust standard errors clustered by firm. *, **, and *** indicate statistical significance at the $10\%$, $5\%$, and $1\%$ levels, respectively.

\end{small}
\end{singlespace}
\vspace{3mm}
\small
  \centering
    \begin{tabular}{lcccccccc}
    \toprule
    \textbf{ } &
      \multicolumn{2}{c}{\textbf{Total Debt}} &  \multicolumn{2}{c}{\textbf{Lev\textsubscript{GAAP}}} &
         \multicolumn{2}{c}{\textbf{PPE}} &  \multicolumn{2}{c}{\textbf{Assets}}
      \\
      \cline{2-9}

              &\multicolumn{1}{c}{(1)}   &\multicolumn{1}{c}{(2)}   &\multicolumn{1}{c}{(3)}   &\multicolumn{1}{c}{(4)}   &\multicolumn{1}{c}{(5)}   &\multicolumn{1}{c}{(6)}   &\multicolumn{1}{c}{(7)}   &\multicolumn{1}{c}{(8)}   \\
\hline
Convert x Post&       0.156   &               &       0.048   &               &      -0.226   &               &      -0.043   &               \\
            &      (0.86)   &               &      (1.31)   &               &     (-1.10)   &               &     (-0.66)   &               \\
Preferred Dummy x Post&               &       0.063   &               &      -0.009   &               &      -0.018   &               &       0.043   \\
            &               &      (1.02)   &               &     (-0.68)   &               &     (-0.39)   &               &      (1.31)   \\
Profitability&      -1.264** &      -0.011   &      -0.412***&       0.037***&      -0.478   &       0.235***&      -0.223   &       0.428***\\
            &     (-2.50)   &     (-0.60)   &     (-3.54)   &      (7.67)   &     (-1.16)   &      (9.46)   &     (-0.62)   &     (19.09)   \\
Tangibility &       0.119   &       0.542***&       0.280** &       0.114***&       2.828***&       3.732***&      -0.126   &      -0.193   \\
            &      (0.34)   &      (3.56)   &      (2.27)   &      (3.22)   &      (5.08)   &     (18.55)   &     (-0.48)   &     (-1.61)   \\
Sales       &       0.398***&       0.370***&       0.021   &       0.021***&       0.426***&       0.346***&       0.369***&       0.338***\\
            &      (5.73)   &      (8.79)   &      (1.46)   &      (3.39)   &      (6.24)   &     (12.45)   &      (8.56)   &     (15.20)   \\
Market-to-Book&      -0.132** &      -0.002***&      -0.019   &      -0.000***&      -0.099***&      -0.004***&      -0.132***&      -0.008***\\
            &     (-2.20)   &     (-4.94)   &     (-0.83)   &     (-3.01)   &     (-2.71)   &     (-8.50)   &     (-3.83)   &    (-18.99)   \\
            &               &               &               &               &               &               &               &               \\
Sample      &        Junk   &     Unrated   &        Junk   &     Unrated   &        Junk   &     Unrated   &        Junk   &     Unrated   \\
Firm FE     &           Y   &           Y   &           Y   &           Y   &           Y   &           Y   &           Y   &           Y   \\
Quarter FE  &           Y   &           Y   &           Y   &               &           Y   &           Y   &           Y   &           Y   \\
Firm Quarters&        4736   &       25858   &        4766   &       35390   &        4766   &       33357   &        4766   &       35388   \\
R\textsuperscript{2}&       0.094   &       0.028   &       0.038   &       0.022   &       0.363   &       0.284   &       0.278   &       0.478   \\


    \bottomrule
    \end{tabular}
\end{table}
\end{landscape}

\clearpage

\newpage


\begin{landscape}
\pagestyle{empty}
\begin{table}[htbp]
\caption{\textbf{Stock Price Response to the Announcement of Moody's Rule Change}
} \label{table_car}
\begin{singlespace}
\begin{small}

\noindent This table contains an event study testing if treated firms experienced positive cumulative abnormal returns compared to matched firms. The dependent variable is the sum of the firm's daily stock return minus the market return over the return period. T-statistics are shown below the parameter estimates in parentheses and are calculated using robust standard errors. *, **, and *** indicate statistical significance at the $10\%$, $5\%$, and $1\%$ levels, respectively.

\end{small}
\end{singlespace}
\vspace{3mm}
\small
  \centering

\begin{tabular}{lcccccccc}
\\
    \toprule
    &
        \multicolumn{2}{c}{\textbf{Pre-Period}} &
      \multicolumn{4}{c}{\textbf{Event}} &
        \multicolumn{2}{c}{\textbf{Post-Period}}
        \\
                  \cline{2-9}
        \\
    &
    \multicolumn{2}{c}{\textbf{CAR[-30,-2]}} &
      \multicolumn{2}{c}{\textbf{CAR[-1,1]}} &
        \multicolumn{2}{c}{\textbf{CAR[-1,3]}} &
           \multicolumn{2}{c}{\textbf{CAR[4,30]}}
      \\
          \cline{2-9}
            &\multicolumn{1}{c}{(1)}   &\multicolumn{1}{c}{(2)}   &\multicolumn{1}{c}{(3)}   &\multicolumn{1}{c}{(4)}   &\multicolumn{1}{c}{(5)}   &\multicolumn{1}{c}{(6)}   &\multicolumn{1}{c}{(7)}   &\multicolumn{1}{c}{(8)}   \\
\hline
Preferred Dummy&       0.009   &               &       0.017** &               &       0.028** &               &       0.002   &               \\
            &      (0.53)   &               &      (2.04)   &               &      (2.60)   &               &      (0.10)   &               \\
Preferred/Capital&               &       0.209   &               &       0.112*  &               &       0.211*  &               &       0.029   \\
            &               &      (1.31)   &               &      (1.72)   &               &      (1.75)   &               &      (0.13)   \\
Profitability&       0.115   &       0.074   &       0.916***&       0.903***&       1.437***&       1.409***&       2.121** &       2.116** \\
            &      (0.18)   &      (0.12)   &      (2.85)   &      (2.79)   &      (3.66)   &      (3.55)   &      (2.27)   &      (2.26)   \\
Tangibility &       0.024   &       0.027   &      -0.002   &      -0.002   &      -0.009   &      -0.010   &      -0.027   &      -0.027   \\
            &      (0.50)   &      (0.56)   &     (-0.06)   &     (-0.08)   &     (-0.29)   &     (-0.30)   &     (-0.43)   &     (-0.43)   \\
Sales       &      -0.007   &      -0.006   &       0.003   &       0.004   &       0.000   &       0.001   &       0.012   &       0.012*  \\
            &     (-1.08)   &     (-0.97)   &      (1.02)   &      (1.18)   &      (0.07)   &      (0.34)   &      (1.65)   &      (1.70)   \\
Market-to-Book&      -0.011   &      -0.005   &       0.022   &       0.020   &       0.027   &       0.023   &      -0.084   &      -0.084   \\
            &     (-0.25)   &     (-0.11)   &      (1.16)   &      (1.03)   &      (1.16)   &      (1.00)   &     (-1.63)   &     (-1.64)   \\
            &               &               &               &               &               &               &               &               \\
Cohort FE   &           Y   &           Y   &           Y   &           Y   &           Y   &           Y   &           Y   &           Y   \\
Firms       &         197   &         197   &         197   &         197   &         197   &         197   &         197   &         197   \\
R\textsuperscript{2}&       0.346   &       0.354   &       0.307   &       0.300   &       0.374   &       0.366   &       0.315   &       0.315   \\

    \bottomrule
    \end{tabular}


\end{table}
\end{landscape}
\clearpage

\begin{landscape}
\pagestyle{empty}
\begin{table}[htbp]
\caption{\textbf{Credit Spread Response to the Announcement of Moody's Rule Change}
} \label{table_spreadsm}
\begin{singlespace}
\begin{small}

\noindent This table contains an event study testing if treated firms' credit spreads increased around the rule change compared to matched firms. The dependent variable is the change in the firm's credit spread over the period in basis points. T-statistics are shown below the parameter estimates in parentheses and are calculated using robust standard errors. *, **, and *** indicate statistical significance at the $10\%$, $5\%$, and $1\%$ levels, respectively.

\end{small}
\end{singlespace}
\vspace{3mm}
\small
  \centering

\begin{tabular}{lcccccccc}
\\
    \toprule
        &
        \multicolumn{2}{c}{\textbf{Pre-Period}} &
      \multicolumn{4}{c}{\textbf{Event}} &
        \multicolumn{2}{c}{\textbf{Post-Period}}
        \\
                  \cline{2-9}
        \\
    &
    \multicolumn{2}{c}{\textbf{$\Delta$ CS[-30,-2]}} &
      \multicolumn{2}{c}{\textbf{$\Delta$ CS[-1,1]}} &
        \multicolumn{2}{c}{\textbf{$\Delta$ CS[-1,3]}} &
           \multicolumn{2}{c}{\textbf{$\Delta$ CS[4,30]}}
      \\
          \cline{2-9}
            &\multicolumn{1}{c}{(1)}   &\multicolumn{1}{c}{(2)}   &\multicolumn{1}{c}{(3)}   &\multicolumn{1}{c}{(4)}   &\multicolumn{1}{c}{(5)}   &\multicolumn{1}{c}{(6)}   &\multicolumn{1}{c}{(7)}   &\multicolumn{1}{c}{(8)}   \\
\hline
Preferred Dummy&       -19.9   &               &        13.5   &               &        16.2*  &               &         2.3   &               \\
            &     (-1.64)   &               &      (1.43)   &               &      (1.70)   &               &      (0.27)   &               \\
Preferred/Capital&               &      -107.9   &               &        74.2*  &               &        96.4** &               &       -52.2   \\
            &               &     (-1.21)   &               &      (1.68)   &               &      (2.26)   &               &     (-1.33)   \\
Profitability&       109.0   &       202.3   &      -489.4*  &      -552.3** &      -385.4   &      -456.1** &        -8.0   &       -60.2   \\
            &      (0.20)   &      (0.38)   &     (-1.97)   &     (-2.36)   &     (-1.62)   &     (-2.12)   &     (-0.03)   &     (-0.22)   \\
Tangibility &        -6.2   &        -0.8   &       -28.4   &       -32.3   &       -15.2   &       -19.9   &        35.3   &        34.4   \\
            &     (-0.18)   &     (-0.02)   &     (-1.44)   &     (-1.49)   &     (-0.82)   &     (-0.96)   &      (1.46)   &      (1.43)   \\
Sales       &        -6.8   &        -7.6   &         3.3   &         3.9   &         3.7   &         4.4   &         5.4   &         5.2   \\
            &     (-1.32)   &     (-1.51)   &      (0.90)   &      (1.03)   &      (1.15)   &      (1.30)   &      (1.27)   &      (1.25)   \\
Market-to-Book&        -5.2   &        -7.1   &        14.0   &        15.1   &        14.7   &        16.1   &         0.7   &         1.0   \\
            &     (-0.24)   &     (-0.34)   &      (0.94)   &      (0.97)   &      (1.02)   &      (1.11)   &      (0.04)   &      (0.06)   \\
            &               &               &               &               &               &               &               &               \\
Cohort FE   &           Y   &           Y   &           Y   &           Y   &           Y   &           Y   &           Y   &           Y   \\
Firms       &          83   &          83   &          88   &          88   &          88   &          88   &          88   &          88   \\
R\textsuperscript{2}&       0.282   &       0.274   &       0.260   &       0.248   &       0.295   &       0.279   &       0.312   &       0.320   \\

    \bottomrule
    \end{tabular}


\end{table}
\end{landscape}
\clearpage
\newpage

\begin{landscape}
\pagestyle{empty}

\begin{table}[htbp]
\caption{\textbf{The Effect of Convertible Debt on Response to Rule Change}} \label{table_convert}
\begin{singlespace}
\begin{small}

\noindent   This table contains results testing whether treated firms' responses to the rule change are affected by the amount of convertible debt they had in their capital structure at the time of the rule change. Preferred Dummy is an indicator variable that equals one if the firm had preferred stock in its capital structure in the last quarter prior to the rule change. Preferred/Capital is the ratio between the amount of preferred stock and the sum of debt and shareholders' equity in the last quarter prior to the rule change. Convert is the ratio of convertible debt to total debt at the time of the rule change. T-statistics are shown below the parameter estimates in parentheses and are calculated using robust standard errors clustered by firm. *, **, and *** indicate statistical significance at the $10\%$, $5\%$, and $1\%$ levels, respectively.
\end{small}
\end{singlespace}
\vspace{3mm}

\small
  \centering
    \begin{tabular}{lcccccccc}
    \toprule
        \textbf{ } &
      \multicolumn{2}{c}{\textbf{Total Debt}} &  \multicolumn{2}{c}{\textbf{Lev\textsubscript{GAAP}}} &
         \multicolumn{2}{c}{\textbf{PPE}} &  \multicolumn{2}{c}{\textbf{Assets}}
      \\
           \cline{2-9}

            &\multicolumn{1}{c}{(1)}   &\multicolumn{1}{c}{(2)}   &\multicolumn{1}{c}{(3)}   &\multicolumn{1}{c}{(4)}   &\multicolumn{1}{c}{(5)}   &\multicolumn{1}{c}{(6)}   &\multicolumn{1}{c}{(7)}   &\multicolumn{1}{c}{(8)}   \\
\hline
Preferred Dummy x Post&       0.290***&               &       0.049***&               &       0.111***&               &       0.117***&               \\
            &      (4.04)   &               &      (2.69)   &               &      (2.77)   &               &      (2.70)   &               \\
Preferred Dummy x Convert x Post&      -1.460***&               &      -0.397*  &               &      -0.657** &               &      -0.734** &               \\
            &     (-3.27)   &               &     (-1.77)   &               &     (-2.00)   &               &     (-2.50)   &               \\
Preferred/Capital x Post&               &       1.445***&               &       0.292** &               &       0.494*  &               &       0.579*  \\
            &               &      (3.65)   &               &      (2.02)   &               &      (1.70)   &               &      (1.73)   \\
Preferred/Capital x Convert x Post&               &      -9.143*  &               &      -3.563** &               &      -1.218   &               &      -0.990   \\
            &               &     (-1.81)   &               &     (-2.44)   &               &     (-0.41)   &               &     (-0.43)   \\
Convert x Post&       0.104   &      -0.013   &      -0.001   &      -0.024   &       0.049   &      -0.027   &       0.055   &      -0.032   \\
            &      (0.48)   &     (-0.06)   &     (-0.01)   &     (-0.38)   &      (0.26)   &     (-0.16)   &      (0.37)   &     (-0.23)   \\
Profitability&      -0.022   &      -0.029   &      -0.223** &      -0.225** &       0.326** &       0.324** &       0.442***&       0.439***\\
            &     (-0.10)   &     (-0.12)   &     (-2.09)   &     (-2.09)   &      (2.45)   &      (2.36)   &      (3.55)   &      (3.41)   \\
Tangibility &       0.171   &       0.187   &       0.360** &       0.362** &       1.948***&       1.967***&      -0.430   &      -0.405   \\
            &      (0.37)   &      (0.40)   &      (2.34)   &      (2.34)   &      (4.57)   &      (4.56)   &     (-1.29)   &     (-1.20)   \\
Sales       &       0.347***&       0.359***&      -0.014   &      -0.012   &       0.347***&       0.353***&       0.334***&       0.340***\\
            &      (5.61)   &      (5.90)   &     (-0.67)   &     (-0.59)   &      (7.18)   &      (7.43)   &      (6.52)   &      (6.74)   \\
Market-to-Book&      -0.051   &      -0.037   &      -0.073** &      -0.070*  &      -0.151** &      -0.145** &      -0.147** &      -0.141** \\
            &     (-0.28)   &     (-0.19)   &     (-2.09)   &     (-1.94)   &     (-2.30)   &     (-2.20)   &     (-2.17)   &     (-2.08)   \\
\_cons      &       4.915***&       4.839***&       0.576***&       0.564***&       3.926***&       3.883***&       6.320***&       6.273***\\
            &     (10.56)   &     (10.34)   &      (4.48)   &      (4.48)   &     (13.30)   &     (13.20)   &     (18.88)   &     (18.80)   \\
            &               &               &               &               &               &               &               &               \\
Firm FE     &           Y   &           Y   &           Y   &           Y   &           Y   &           Y   &           Y   &           Y   \\
Quarter x Cohort FE&           Y   &           Y   &           Y   &           Y   &           Y   &           Y   &           Y   &           Y   \\
Firm Quarters&        2448   &        2448   &        2454   &        2454   &        2454   &        2454   &        2454   &        2454   \\
R\textsuperscript{2}&       0.139   &       0.118   &       0.062   &       0.058   &       0.342   &       0.334   &       0.309   &       0.299   \\



    \bottomrule
    \end{tabular}
\end{table}
\end{landscape}


\clearpage
\newpage

\begin{table}[htbp]
\caption{\textbf{The Effect of Rule Change on Preferred Stock Levels}
} \label{table_issue}
\begin{singlespace}
\begin{small}

\noindent This table contains results testing whether firms only rated speculative-grade by Moody's at the time of the rule change increased their preferred stock levels relative to other firms rated speculative grade by either Moody's or S\&P at the time of the rule change. The dependent variables are preferred divided by capital in percentage points (P/C) and the log of total preferred (Log(P)). MJOnly is a dummy variable that equals 1 if the firm is only rated speculative-grade by Moody's. MIGOnly is a dummy variable that equals 1 if the firm is only rated IG by Moody's. The sample is restricted to firms that are rated speculative-grade by either Moody's or S\&P at the time of the rule change in Columns (1) and (2). In Columns (3) and (4), the sample is restricted to firms that are rated investment-grade by either Moody's or S\&P at the time of the rule change. T-statistics are shown below the parameter estimates in parentheses and are calculated using robust standard errors clustered by firm. *, **, and *** indicate statistical significance at the $10\%$, $5\%$, and $1\%$ levels, respectively.

\end{small}
\end{singlespace}
\vspace{3mm}
\small
  \centering

\begin{tabular}{lcc|cc}
    \toprule &  \multicolumn{2}{c|}{\textbf{Main Test}} &  \multicolumn{2}{c}{\textbf{Placebo}}
    \\
     &  \multicolumn{1}{c}{\textbf{P/C }} &
       \multicolumn{1}{c|}{\textbf{Log(P)}} &
  \multicolumn{1}{c}{\textbf{P/C }} &
       \multicolumn{1}{c}{\textbf{Log(P)}}
      \\
          \cline{2-5}
            &\multicolumn{1}{c}{(1)}   &\multicolumn{1}{c}{(2)}   &\multicolumn{1}{c}{(3)}   &\multicolumn{1}{c}{(4)}   \\
\hline
MJOnly x Post&       0.543** &       0.158** &               &               \\
            &      (2.36)   &      (2.47)   &               &               \\
MIGOnly x Post&               &               &      -0.333   &      -0.352   \\
            &               &               &     (-1.08)   &     (-1.09)   \\
Profitability&      -1.081   &      -0.400   &       7.886   &       2.234   \\
            &     (-1.23)   &     (-1.32)   &      (1.13)   &      (1.36)   \\
Tangibility &       1.907   &      -0.603   &     -10.453   &       0.617   \\
            &      (0.93)   &     (-0.88)   &     (-0.98)   &      (1.05)   \\
Sales       &       0.254   &       0.200*  &      -1.091   &      -0.016   \\
            &      (1.55)   &      (1.95)   &     (-1.04)   &     (-0.18)   \\
Market-to-Book&      -0.184   &      -0.009   &       0.275   &       0.131** \\
            &     (-1.32)   &     (-0.15)   &      (1.27)   &      (2.40)   \\
Constant    &      -1.512   &      -0.668   &      11.350   &       0.015   \\
            &     (-1.17)   &     (-1.06)   &      (1.04)   &      (0.02)   \\
            &               &               &               &               \\
Sample      &        Junk   &        Junk   &          IG   &          IG   \\
Firm FE     &           Y   &           Y   &           Y   &           Y   \\
Quarter FE  &           Y   &           Y   &           Y   &           Y   \\
Firm Quarters&        6285   &        6285   &        5074   &        5074   \\
R\textsuperscript{2}&       0.003   &       0.008   &       0.081   &       0.013   \\


    \bottomrule
    \end{tabular}


\end{table}


\clearpage

\newpage


\begin{table}[htbp]
\caption{\textbf{The Effect of Rule Change on Credit Ratings}} \label{table_cr}
\begin{singlespace}
\begin{small}

\noindent   This table contains results testing whether treated firms' ratings change after the rule change. The dependent variables, Moody's Rating and S\&P Rating, are categorical variables that take values between 1 and 22 that are mapped from Moody's and S\&P letter ratings, where 1 is the highest rating and 22 the lowest. Preferred Dummy is an indicator variable that equals one if the firm had preferred stock in its capital structure in the last quarter prior to the rule change. Preferred/Capital is the ratio between the amount of preferred stock and the sum of debt and shareholders' equity in the last quarter prior to the rule change.  T-statistics are shown below the parameter estimates in parentheses and are calculated using robust standard errors clustered by firm. *, **, and *** indicate statistical significance at the $10\%$, $5\%$, and $1\%$ levels, respectively.
\end{small}
\end{singlespace}
\vspace{3mm}

\small
  \centering
    \begin{tabular}{lcccc}
    \toprule
    \textbf{ } &
      \multicolumn{2}{c}{\textbf{Moody's Rating}} &  \multicolumn{2}{c}{\textbf{S\&P Rating}}

      \\
           \cline{2-5}

            &\multicolumn{1}{c}{(1)}   &\multicolumn{1}{c}{(2)}   &\multicolumn{1}{c}{(3)}   &\multicolumn{1}{c}{(4)}   \\
\hline
Preferred Dummy x Post&      -0.103   &               &      -0.219   &               \\
            &     (-0.57)   &               &     (-1.56)   &               \\
Preferred/Capital x Post&               &      -1.423   &               &      -1.000   \\
            &               &     (-1.17)   &               &     (-1.08)   \\
Profitability&      -1.135   &      -1.125   &      -2.752***&      -2.753***\\
            &     (-1.19)   &     (-1.17)   &     (-2.83)   &     (-2.84)   \\
Tangibility &       2.271*  &       2.227   &       2.145   &       2.122   \\
            &      (1.68)   &      (1.64)   &      (1.63)   &      (1.61)   \\
Sales       &      -0.178   &      -0.177   &      -0.776***&      -0.782***\\
            &     (-0.70)   &     (-0.70)   &     (-4.32)   &     (-4.35)   \\
Market-to-Book&      -1.145***&      -1.132***&      -1.053***&      -1.043***\\
            &     (-3.39)   &     (-3.34)   &     (-3.35)   &     (-3.33)   \\
\_cons      &      16.019***&      16.026***&      18.635***&      18.653***\\
            &     (11.98)   &     (11.97)   &     (17.08)   &     (17.16)   \\
            &               &               &               &               \\
Firm FE     &           Y   &           Y   &           Y   &           Y   \\
Quarter x Cohort FE&           Y   &           Y   &           Y   &           Y   \\
Firm Quarters&        2340   &        2340   &        2222   &        2222   \\
R\textsuperscript{2}&       0.079   &       0.081   &       0.212   &       0.208   \\



    \bottomrule
    \end{tabular}
\end{table}


\clearpage
\newpage


\newpage


\section*{Online Appendix}
\renewcommand{\thefigure}{A.\arabic{figure}}
\renewcommand{\thetable}{A.\arabic{table}}
\setcounter{figure}{0}
\setcounter{table}{0}


\begin{table}[htbp]
\caption{\textbf{List of Treated Firms}} \label{table_list}
\begin{singlespace}
\begin{small}

\end{small}
\end{singlespace}
\vspace{3mm}

\small
  \centering
    \begin{tabular}{ll}
    \toprule


Company Name (Compustat) & Fama French 49 Industry \\
\hline
Alcoa Inc&Steel Works Etc\\
Alere Inc&Pharmaceutical Products\\
Atlas Pipeline Partner LP&Utilities\\
Century Aluminum Co&Steel Works Etc\\
Chesapeake Energy Corp&Petroleum and Natural Gas\\
Cincinnati Bell Inc&Communication\\
Colfax Corp&Machinery\\
Cooper-Standard Holdings Inc&Automobiles and Trucks\\
Cumulus Media Inc&Communication\\
Dana Inc&Automobiles and Trucks\\
Endeavour International Corp&Petroleum and Natural Gas\\
EnLink Midstream Partners LP&Utilities\\
Erickson Inc&Business Services\\
Forbes Energy Services Ltd&Petroleum and Natural Gas\\
Gastar Exploration Inc&Petroleum and Natural Gas\\
General Cable Corp/De&Steel Works Etc\\
General Motors Co&Automobiles and Trucks\\
Goodrich Petroleum Corp&Petroleum and Natural Gas\\
HealthSouth Corp&Healthcare\\
Hecla Mining Co&Precious Metals\\
Hovnanian Entrprs Inc  -Cl A&Construction\\
HRG Group Inc&Electrical Equipment\\
ION Geophysical Corp&Measuring and Control Equipment\\
LSB Industries Inc&Chemicals\\
M/I Homes Inc&Construction\\
Magnum Hunter Resources Corp&Petroleum and Natural Gas\\
Navistar International Corp&Automobiles and Trucks\\
NRG Energy Inc&Utilities\\
Nuance Communications Inc&Computer Software\\
Office Depot Inc&Retail\\
Officemax Inc&Wholesale\\
Penn Virginia Corp&Petroleum and Natural Gas\\
PetroQuest Energy Inc&Petroleum and Natural Gas\\
PNM Resources Inc&Utilities\\
Post Holdings Inc&Food Products\\
Regency Energy Partners Lp&Petroleum and Natural Gas\\
Rite Aid Corp&Retail\\
Sanchez Energy Corp&Petroleum and Natural Gas\\
SandRidge Energy Inc&Petroleum and Natural Gas\\
Spanish Broadcasting Sys Inc&Communication\\
Universal Corp/Va&Wholesale\\
Vanguard Natural Resources&Petroleum and Natural Gas\\
Warren Resources Inc&Petroleum and Natural Gas\\
Westmoreland Coal Co&Coal\\

    \bottomrule
    \end{tabular}
\end{table}

\clearpage
\newpage


\begin{figure}[bp!]
\begin{center}
\begin{minipage}{0.32\textwidth}
    \includegraphics[width=\textwidth]{RCFS Resubmission R2/KZ_w_density.png}
\end{minipage}
\begin{minipage}{0.32\textwidth}
    \includegraphics[width=\textwidth]{RCFS Resubmission R2/WW_w_density.png}
\end{minipage}
\begin{minipage}{0.32\textwidth}
    \includegraphics[width=\textwidth]{RCFS Resubmission R2/SA_w_density.png}
\end{minipage}
\end{center}
\begin{singlespace}
\begin{small}
\caption{\textbf{Financial Constraints Histogram}}
\label{fig_fincon}
\noindent The figure includes a histogram of the Kaplan-Zingales, Whited-Wu, and Size-Age indices of financial constraints for the firms in the sample, grouped by treated, matched, untreated, and all other speculative firms.
\end{small}
\end{singlespace}
\end{figure}


%%% Table created in Stata by command iebaltab
%%% (https://github.com/worldbank/ietoolkit)
%%% (https://dimewiki.worldbank.org/iebaltab)
%%% The command was specified exactly like this: 
%%% iebaltab WW_w KZ_w SA_w if TMJunk == 1, grpvar(treated_matched) grplabels(2 "Treated" @ 1 "Matchd Untreated" @ 0 "Unmatched Untreated") rowvarlabels rowlabels(WW_w "Whited-Wu Index" KZ_w "Kaplan-Zingales Index" SA_w "Size-Age Index") savetex("C:/Users/cf8745/Box/Research/Rating Agency Paper - Final/ReplicationPackage/Output/Tables/fin_constraints_balance.tex") replace addnote("* p<0.10, ** p<0.05, *** p<0.01")

\begin{tabular}{@{\extracolsep{5pt}}lcccccccccccc}
\\[-1.8ex]\hline \hline \\[-1.8ex]
 & \multicolumn{2}{c}{(1)}  & \multicolumn{2}{c}{(2)}  & \multicolumn{2}{c}{(3)}  & \multicolumn{2}{c}{(1)-(2)} & \multicolumn{2}{c}{(1)-(3)} & \multicolumn{2}{c}{(2)-(3)} \\
 & \multicolumn{2}{c}{Unmatched Untreated}  & \multicolumn{2}{c}{Matchd Untreated}  & \multicolumn{2}{c}{Treated}  & \multicolumn{6}{c}{Pairwise t-test}  \\
Variable & N & Mean/(SE) & N & Mean/(SE) & N & Mean/(SE) & N & Mean difference & N & Mean difference & N & Mean difference \\ \hline \\[-1.8ex] 
Whited-Wu Index KZ\_w Kaplan-Zingales Index SA\_w Size-Age Index   & 315    & -0.310    & 108    & -0.317    & 43    & -0.326    & 423    & 0.006    & 358    & 0.016    & 151    & 0.009   \\
 &   & (0.006)  &   & (0.008)  &   & (0.015)  &   &  &   &  &   &  \\ [1ex]
Kaplan-Zingales Index (winsorized)   & 303    & 1.276    & 108    & 0.001    & 43    & 1.002    & 411    & 1.275**    & 346    & 0.274    & 151    & -1.001   \\
 &   & (0.156)  &   & (0.825)  &   & (0.655)  &   &  &   &  &   &  \\ [1ex]
Size-Age Index (winsorized)   & 317    & -3.369    & 108    & -3.390    & 44    & -3.369    & 425    & 0.021    & 361    & -0.000    & 152    & -0.021   \\
 &   & (0.007)  &   & (0.007)  &   & (0.016)  &   &  &   &  &   &  \\ [1ex]
\hline \hline \\[-1.8ex]
%%% This is the note. If it does not have the correct margins, use texnotewidth() option or change the number before '\textwidth' in line below to fit it to table size.
\multicolumn{10}{@{} p{\textwidth}}{If the table includes missing values (.n, .o, .v etc.) see the Missing values section in the help file for the Stata command iebaltab for definitions of these values. Significance: ***=.01, **=.05, *=.1. Full user input as written by user: [iebaltab WW\_w KZ\_w SA\_w if TMJunk == 1, grpvar(treated\_matched) grplabels(2 "Treated" @ 1 "Matchd Untreated" @ 0 "Unmatched Untreated") rowvarlabels rowlabels(WW\_w "Whited-Wu Index" KZ\_w "Kaplan-Zingales Index" SA\_w "Size-Age Index") savetex("C:/Users/cf8745/Box/Research/Rating Agency Paper - Final/ReplicationPackage/Output/Tables/fin\_constraints\_balance.tex") replace addnote("* p<0.10, ** p<0.05, *** p<0.01")] * p<0.10, ** p<0.05, *** p<0.01}

\end{tabular}



\clearpage

\newpage

\begin{table}[htbp]
\caption{\textbf{Ratings and Lev\textsubscript{GAAP}  }
} \label{table_ratbook}
\begin{singlespace}
\begin{small}

\noindent The top table displays the average $Lev_{GAAP}$ levels for speculative-grade firms in July 2013. Moody's Rating is a categorical variable that can take values between 1 and 22, which are mapped from Moody's letter ratings, where 1 is the highest rating and 22 the lowest. The bottom portion includes a simple cross-sectional regression in July of 2013 with $Lev_{GAAP}$ as the dependent variable and the firm's Moody's rating as the independent variable. In Column (1), the sample is all speculative-grade firms, while in Column (2), the sample is firms rated between Caa1 and Ba3.

\end{small}
\end{singlespace}
\vspace{3mm}
\small
  \centering
{
\def\sym#1{\ifmmode^{#1}\else\(^{#1}\)\fi}
\begin{tabular}{l*{3}{ccccccc}}               \\
\hline
                    &\multicolumn{1}{c}{(1)}&\multicolumn{1}{c}{(2)}&\multicolumn{1}{c}{(3)}&\multicolumn{1}{c}{(4)}&\multicolumn{1}{c}{(5)}&\multicolumn{1}{c}{(6)}&\multicolumn{1}{c}{(7)}&\multicolumn{1}{c}{(8)}&\multicolumn{1}{c}{(9)}&\multicolumn{1}{c}{(10)}\\
\hline
Book Leverage       &       0.430&       0.463&       0.498&       0.590&       0.659&       0.632&       0.730&       0.858&       0.356&       0.942\\
N                   &          41&          51&          85&          86&          72&          77&          41&          18&           2&           2\\


\hline


\end{tabular}
}           \\


\begin{tabular}{lcc}
\\
    \toprule
    \textbf{} &
      \multicolumn{2}{c}{\textbf{Lev\textsubscript{GAAP}}}
      \\
          \cline{2-3}
    \textbf{} &
      \textbf{(1)} &
      \textbf{(2)}
      \\
       \midrule
                &\multicolumn{1}{c}{(1)}   &\multicolumn{1}{c}{(2)}   \\
\hline
Moody's Rating&       0.053***&       0.060***\\
            &      (9.75)   &      (6.72)   \\
Sample      &         All   &  Caa1 - Ba3   \\
Industry FE &           Y   &           Y   \\
N           &         475   &         357   \\
R\textsuperscript{2}&       0.346   &       0.321   \\

    \bottomrule
    \end{tabular}


\end{table}
\clearpage
\newpage

\begin{table}[htbp]
\caption{\textbf{The Effect of S\&P Rating on Response to Rule Change}
} \label{table_cross}
\begin{singlespace}
\begin{small}

\noindent This table contains results testing whether treated firms respond to the rule change less if they are rated by S\&P.  The dependent variables are Total Debt, PPE, Assets and Moody's Rating. Preferred Dummy is an indicator variable that equals one if the firm had preferred stock in its capital structure in the last quarter prior to the rule change. Preferred/Capital is the ratio between the amount of preferred stock and the sum of debt and shareholders' equity in the last quarter prior to the rule change. SP is an indicator variable that equals 1 if the firm is rated by S\&P at the time of the rule change. T-statistics are shown below the parameter estimates in parentheses and are calculated using robust standard errors clustered by firm. *, **, and *** indicate statistical significance at the $10\%$, $5\%$, and $1\%$ levels, respectively.

\end{small}
\end{singlespace}
\vspace{3mm}
\small
  \centering
   \scalebox{1}{ \begin{tabular}{lcccc}
    \toprule
    \textbf{ } &
      \multicolumn{1}{c}{\textbf{Total Debt}} &
       \multicolumn{1}{c}{\textbf{PPE}} &
       \multicolumn{1}{c}{\textbf{Assets}} &
       \multicolumn{1}{c}{\textbf{Moody's Rating}}
      \\
          \cline{2-5}
              &\multicolumn{1}{c}{(1)}   &\multicolumn{1}{c}{(2)}   &\multicolumn{1}{c}{(3)}   &\multicolumn{1}{c}{(4)}   \\
\hline
Preferred Dummy x Post x SP&      -0.285   &      -0.201   &      -0.284** &      -0.614   \\
            &     (-1.28)   &     (-1.57)   &     (-2.15)   &     (-1.41)   \\
Preferred Dummy x Post&       0.485** &       0.267** &       0.344***&       0.453   \\
            &      (2.33)   &      (2.27)   &      (2.83)   &      (1.24)   \\
Post x SP   &       0.123   &       0.127** &       0.129*  &       0.014   \\
            &      (1.32)   &      (1.99)   &      (1.79)   &      (0.05)   \\
Profitability&       0.000   &       0.340** &       0.458***&      -1.135   \\
            &      (0.00)   &      (2.54)   &      (3.66)   &     (-1.19)   \\
Tangibility &       0.217   &       1.968***&      -0.403   &       2.314*  \\
            &      (0.47)   &      (4.66)   &     (-1.22)   &      (1.71)   \\
Sales       &       0.356***&       0.352***&       0.339***&      -0.187   \\
            &      (5.91)   &      (7.50)   &      (6.86)   &     (-0.74)   \\
Market-to-Book&      -0.038   &      -0.146** &      -0.142** &      -1.155***\\
            &     (-0.21)   &     (-2.25)   &     (-2.11)   &     (-3.42)   \\
\_cons      &       4.752***&       3.806***&       6.197***&      16.057***\\
            &     (10.40)   &     (13.08)   &     (19.39)   &     (11.88)   \\
            &               &               &               &               \\
Firm FE     &           Y   &           Y   &           Y   &           Y   \\
Quarter x Cohort FE&           Y   &           Y   &           Y   &           Y   \\
Firm Quarters&        2448   &        2454   &        2454   &        2340   \\
R\textsuperscript{2}&       0.131   &       0.342   &       0.310   &       0.082   \\



    \bottomrule
    \end{tabular}}
\end{table}


\clearpage

\newpage

\subsection*{Unmatched Regressions}

\begin{table}[htbp]
\caption{\textbf{The Effect of Rule Change on Debt Levels}} \label{table_debt_all}
\begin{singlespace}
\begin{small}

\noindent This table contains results testing whether treated firms increase their debt levels after the rule change compared to other firms rated speculative-grade by Moody's. The dependent variables are total debt, long-term debt and short-term debt and are all in logs. Preferred Dummy is an indicator variable that equals one if the firm had preferred stock in its capital structure in the last quarter prior to the rule change. Preferred/Capital is the ratio between the amount of preferred stock and the sum of debt and shareholders' equity in the last quarter prior to the rule change.  T-statistics are shown below the parameter estimates in parentheses and are calculated using robust standard errors clustered by firm. *, **, and *** indicate statistical significance at the $10\%$, $5\%$, and $1\%$ levels, respectively.

\end{small}
\end{singlespace}
\vspace{3mm}
\small
  \centering
    \begin{tabular}{lcccccc}
    \toprule
    \textbf{} &
      \multicolumn{2}{c}{\textbf{Total Debt}} &  \multicolumn{2}{c}{\textbf{LT Debt}} &
      \multicolumn{2}{c}{\textbf{ST Debt}}
      \\
      \cline{2-7}

            &\multicolumn{1}{c}{(1)}   &\multicolumn{1}{c}{(2)}   &\multicolumn{1}{c}{(3)}   &\multicolumn{1}{c}{(4)}   &\multicolumn{1}{c}{(5)}   &\multicolumn{1}{c}{(6)}   \\
\hline
Preferred Dummy x Post&       0.204***&               &       0.236***&               &       0.059   &               \\
            &      (2.68)   &               &      (2.75)   &               &      (0.30)   &               \\
Preferred/Capital x Post&               &       1.547***&               &       1.662***&               &       0.126   \\
            &               &      (3.71)   &               &      (3.36)   &               &      (0.05)   \\
Profitability&      -1.085***&      -1.089***&      -2.241***&      -2.234***&      -2.776***&      -2.784***\\
            &     (-2.62)   &     (-2.66)   &     (-5.25)   &     (-5.31)   &     (-3.96)   &     (-3.98)   \\
Tangibility &       0.318   &       0.331   &       0.269   &       0.285   &      -0.537   &      -0.537   \\
            &      (1.01)   &      (1.04)   &      (0.76)   &      (0.80)   &     (-0.62)   &     (-0.62)   \\
Sales       &       0.414***&       0.411***&       0.423***&       0.420***&       0.316** &       0.317** \\
            &      (6.52)   &      (6.53)   &      (4.86)   &      (4.84)   &      (2.02)   &      (2.03)   \\
Market-to-Book&      -0.141** &      -0.145** &      -0.130** &      -0.135** &      -0.176   &      -0.178   \\
            &     (-2.42)   &     (-2.49)   &     (-2.29)   &     (-2.38)   &     (-1.23)   &     (-1.25)   \\
\_cons      &       4.468***&       4.488***&       4.365***&       4.386***&       1.605   &       1.603   \\
            &     (10.82)   &     (10.98)   &      (7.80)   &      (7.87)   &      (1.57)   &      (1.57)   \\
            &               &               &               &               &               &               \\
Firm FE     &           Y   &           Y   &           Y   &           Y   &           Y   &           Y   \\
Quarter FE  &           Y   &           Y   &           Y   &           Y   &           Y   &           Y   \\
Firm Quarters&        5259   &        5259   &        5228   &        5228   &        4133   &        4133   \\
R\textsuperscript{2}&       0.110   &       0.111   &       0.095   &       0.094   &       0.009   &       0.009   \\



    \bottomrule
    \end{tabular}
\end{table}
\clearpage

\newpage


\begin{landscape}
\pagestyle{empty}
\begin{table}[htbp]
\caption{\textbf{Yearly Placebo Tests}} \label{table_placebo_all}
\begin{singlespace}
\begin{small}

This table contains results testing whether there are similar changes in total debt in the second half of other years in which the rule change did not take place among treated and untreated firms rated speculative-grade by Moody's. Preferred Dummy is an indicator variable that equals one if the firm had preferred stock in its capital structure in the last quarter prior to the rule change. After July is a dummy variable that equals 1 if the quarter end is after July. Controls include: profitability, tangibility, sales and market-to-book. T-statistics are shown below the parameter estimates in parentheses and are calculated using robust standard errors clustered by firm. *, **, and *** indicate statistical significance at the $10\%$, $5\%$, and $1\%$ levels, respectively
\end{small}
\end{singlespace}
\vspace{3mm}

\small
  \centering
\begin{tabular}{lcccccccc}
    \toprule
    \textbf{} &
      \multicolumn{8}{c}{\textbf{Total Debt}}
      \\
                \cline{2-9}

            &\multicolumn{1}{c}{(1)}   &\multicolumn{1}{c}{(2)}   &\multicolumn{1}{c}{(3)}   &\multicolumn{1}{c}{(4)}   &\multicolumn{1}{c}{(5)}   &\multicolumn{1}{c}{(6)}   &\multicolumn{1}{c}{(7)}   &\multicolumn{1}{c}{(8)}   \\
\hline
Preferred x After July&       0.041   &      -0.023   &      -0.125***&      -0.025   &      -0.039   &       0.173***&       0.049   &       0.004   \\
            &      (0.54)   &     (-0.43)   &     (-2.66)   &     (-0.59)   &     (-1.02)   &      (3.04)   &      (1.16)   &      (0.14)   \\
Constant    &       6.165***&       5.937***&       4.651***&       4.379***&       5.779***&       4.920***&       7.217***&       5.327***\\
            &     (21.13)   &     (11.64)   &      (3.78)   &      (4.16)   &      (9.21)   &      (6.74)   &     (17.03)   &      (6.48)   \\
            &               &               &               &               &               &               &               &               \\
Year        &        2008   &        2009   &        2010   &        2011   &        2012   &        2013   &        2014   &        2015   \\
Controls    &           Y   &           Y   &           Y   &           Y   &           Y   &           Y   &           Y   &           Y   \\
Firm FE     &           Y   &           Y   &           Y   &           Y   &           Y   &           Y   &           Y   &           Y   \\
Quarter FE  &           Y   &           Y   &           Y   &           Y   &           Y   &           Y   &           Y   &           Y   \\
Firm Quarters&        1533   &        1556   &        1615   &        1691   &        1758   &        1791   &        1729   &        1639   \\
R\textsuperscript{2}&       0.026   &       0.005   &       0.045   &       0.107   &       0.116   &       0.081   &       0.047   &       0.032   \\



    \bottomrule
    \end{tabular}
\end{table}
\end{landscape}
\clearpage
\newpage


\begin{table}[htbp]
\caption{\textbf{The Effect of Rule Change on Leverage}
} \label{table_lev_all}
\begin{singlespace}
\begin{small}

\noindent This table contains results testing whether treated firms increase their leverage after the rule change compared to other firms rated speculative-grade by Moody's. The dependent variables are $Lev_{GAAP}$ and $Lev_{Moody's}$. Preferred Dummy is an indicator variable that equals one if the firm had preferred stock in its capital structure in the last quarter prior to the rule change. Preferred/Capital is the ratio between the amount of preferred stock and the sum of debt and shareholders' equity in the last quarter prior to the rule change. T-statistics are shown below the parameter estimates in parentheses and are calculated using robust standard errors clustered by firm. *, **, and *** indicate statistical significance at the $10\%$, $5\%$, and $1\%$ levels, respectively.

\end{small}
\end{singlespace}
\vspace{3mm}
\small
  \centering
   \scalebox{1}{ \begin{tabular}{lcccc}
    \toprule
    \textbf{ } &
      \multicolumn{2}{c}{\textbf{Lev\textsubscript{GAAP}}} &
      \multicolumn{2}{c}{\textbf{Lev\textsubscript{Moody's}}}
      \\
          \cline{2-5}


            &\multicolumn{1}{c}{(1)}   &\multicolumn{1}{c}{(2)}   &\multicolumn{1}{c}{(3)}   &\multicolumn{1}{c}{(4)}   \\
\hline
Preferred Dummy x Post&       0.043** &               &      -0.010   &               \\
            &      (2.36)   &               &     (-0.61)   &               \\
Preferred/Capital x Post&               &       0.471***&               &      -0.148   \\
            &               &      (3.61)   &               &     (-1.44)   \\
Profitability&      -0.325*  &      -0.320*  &      -0.310*  &      -0.312*  \\
            &     (-1.92)   &     (-1.92)   &     (-1.91)   &     (-1.90)   \\
Tangibility &       0.303***&       0.305***&       0.296***&       0.296***\\
            &      (2.69)   &      (2.72)   &      (2.63)   &      (2.63)   \\
Sales       &       0.019   &       0.016   &       0.011   &       0.012   \\
            &      (1.22)   &      (1.12)   &      (0.78)   &      (0.84)   \\
Market-to-Book&      -0.022   &      -0.023   &      -0.023   &      -0.023   \\
            &     (-1.00)   &     (-1.04)   &     (-1.04)   &     (-1.03)   \\
\_cons      &       0.377***&       0.394***&       0.438***&       0.432***\\
            &      (3.25)   &      (3.61)   &      (4.17)   &      (4.08)   \\
            &               &               &               &               \\
Firm FE     &           Y   &           Y   &           Y   &           Y   \\
Quarter FE  &           Y   &           Y   &           Y   &           Y   \\
Firm Quarters&        5291   &        5291   &        5291   &        5291   \\
R\textsuperscript{2}&       0.040   &       0.046   &       0.032   &       0.033   \\


    \bottomrule
    \end{tabular}}
\end{table}


\clearpage

\newpage

\begin{table}[htbp]
\caption{\textbf{The Effect of Rule Change on the Balance Sheet}} \label{table_bs_all}
\begin{singlespace}
\begin{small}

\noindent This table contains results testing whether treated firms increase the size of their balance sheets after the rule change compared with other firms rated speculative-grade by Moody's. The dependent variables are Capex, Assets, and PPE, all of which are in logs. Preferred Dummy is an indicator variable that equals one if the firm had preferred stock in its capital structure in the last quarter prior to the rule change. Preferred/Capital is the ratio between the amount of preferred stock and the sum of debt and shareholders' equity in the last quarter prior to the rule change.  T-statistics are shown below the parameter estimates in parentheses and are calculated using robust standard errors clustered by firm. *, **, and *** indicate statistical significance at the $10\%$, $5\%$, and $1\%$ levels, respectively.

\end{small}
\end{singlespace}
\vspace{3mm}
\small
  \centering
    \begin{tabular}{lcccccc}
    \toprule
    \textbf{ } &
      \multicolumn{2}{c}{\textbf{Capex}} &  \multicolumn{2}{c}{\textbf{Assets}} &
      \multicolumn{2}{c}{\textbf{PPE}}
      \\
      \cline{2-7}


            &\multicolumn{1}{c}{(1)}   &\multicolumn{1}{c}{(2)}   &\multicolumn{1}{c}{(3)}   &\multicolumn{1}{c}{(4)}   &\multicolumn{1}{c}{(5)}   &\multicolumn{1}{c}{(6)}   \\
\hline
Preferred Dummy x Post&       0.046   &               &       0.053   &               &       0.048   &               \\
            &      (0.72)   &               &      (1.43)   &               &      (1.30)   &               \\
Preferred/Capital x Post&               &       0.366   &               &       0.584** &               &       0.428*  \\
            &               &      (1.08)   &               &      (1.97)   &               &      (1.70)   \\
Profitability&      -0.083   &      -0.081   &      -0.116   &      -0.109   &      -0.364   &      -0.360   \\
            &     (-0.26)   &     (-0.25)   &     (-0.42)   &     (-0.40)   &     (-1.08)   &     (-1.07)   \\
Tangibility &       1.060***&       1.063***&      -0.053   &      -0.050   &       2.785***&       2.787***\\
            &      (2.94)   &      (2.94)   &     (-0.23)   &     (-0.22)   &      (5.67)   &      (5.67)   \\
Sales       &       0.369***&       0.368***&       0.389***&       0.386***&       0.434***&       0.432***\\
            &      (6.88)   &      (6.81)   &      (9.90)   &      (9.94)   &      (7.16)   &      (7.10)   \\
Market-to-Book&       0.151***&       0.150***&      -0.135***&      -0.136***&      -0.105***&      -0.106***\\
            &      (3.14)   &      (3.14)   &     (-4.07)   &     (-4.11)   &     (-3.08)   &     (-3.11)   \\
\_cons      &       0.553   &       0.563   &       5.715***&       5.736***&       2.945***&       2.958***\\
            &      (1.57)   &      (1.59)   &     (21.40)   &     (21.77)   &      (5.60)   &      (5.61)   \\
            &               &               &               &               &               &               \\
Firm FE     &           Y   &           Y   &           Y   &           Y   &           Y   &           Y   \\
Quarter FE  &           Y   &           Y   &           Y   &           Y   &           Y   &           Y   \\
Firm Quarters&        5278   &        5278   &        5291   &        5291   &        5291   &        5291   \\
R\textsuperscript{2}&       0.054   &       0.054   &       0.300   &       0.303   &       0.373   &       0.374   \\



    \bottomrule
    \end{tabular}
\end{table}
\clearpage

\newpage

\begin{landscape}
\pagestyle{empty}
\begin{table}[htbp]
\caption{\textbf{Stock Price Response to the Announcement of Moody's Rule Change}
} \label{table_carun}
\begin{singlespace}
\begin{small}

\noindent This table contains an event study testing if treated firms experienced positive cumulative abnormal returns compared to other firms rated speculative-grade by Moody's. The dependent variable is the sum of the firm's daily stock return minus the market return over the return period. T-statistics are shown below the parameter estimates in parentheses and are calculated using robust standard errors. *, **, and *** indicate statistical significance at the $10\%$, $5\%$, and $1\%$ levels, respectively.

\end{small}
\end{singlespace}
\vspace{3mm}
\small
  \centering

\begin{tabular}{lcccccccc}
\\
    \toprule
    &
        \multicolumn{2}{c}{\textbf{Pre-Period}} &
      \multicolumn{4}{c}{\textbf{Event}} &
        \multicolumn{2}{c}{\textbf{Post-Period}}
        \\
                  \cline{2-9}
        \\
    &
    \multicolumn{2}{c}{\textbf{CAR[-30,-2]}} &
      \multicolumn{2}{c}{\textbf{CAR[-1,1]}} &
        \multicolumn{2}{c}{\textbf{CAR[-1,3]}} &
           \multicolumn{2}{c}{\textbf{CAR[4,30]}}
      \\
          \cline{2-9}
            &\multicolumn{1}{c}{(1)}   &\multicolumn{1}{c}{(2)}   &\multicolumn{1}{c}{(3)}   &\multicolumn{1}{c}{(4)}   &\multicolumn{1}{c}{(5)}   &\multicolumn{1}{c}{(6)}   &\multicolumn{1}{c}{(7)}   &\multicolumn{1}{c}{(8)}   \\
\hline
Preferred Dummy&      -0.018   &               &       0.014   &               &       0.028** &               &       0.013   &               \\
            &     (-0.78)   &               &      (1.43)   &               &      (2.29)   &               &      (0.63)   &               \\
Preferred/Capital&               &      -0.008   &               &       0.078   &               &       0.170   &               &       0.148   \\
            &               &     (-0.05)   &               &      (1.21)   &               &      (1.62)   &               &      (0.82)   \\
Profitability&      -1.657   &      -1.647   &       0.816** &       0.803** &       0.883** &       0.855** &       0.671   &       0.653   \\
            &     (-0.99)   &     (-0.99)   &      (2.56)   &      (2.52)   &      (2.36)   &      (2.29)   &      (1.17)   &      (1.15)   \\
Tangibility &      -0.068   &      -0.066   &      -0.000   &      -0.001   &      -0.019   &      -0.020   &       0.024   &       0.024   \\
            &     (-0.86)   &     (-0.84)   &     (-0.02)   &     (-0.05)   &     (-0.81)   &     (-0.87)   &      (0.70)   &      (0.69)   \\
Sales       &      -0.001   &      -0.001   &       0.003   &       0.003   &       0.003   &       0.003   &       0.012*  &       0.013** \\
            &     (-0.15)   &     (-0.13)   &      (0.93)   &      (0.98)   &      (0.67)   &      (0.81)   &      (1.88)   &      (1.99)   \\
Market-to-Book&      -0.016   &      -0.014   &       0.000   &       0.000   &      -0.002   &      -0.002   &      -0.005   &      -0.004   \\
            &     (-0.95)   &     (-0.85)   &      (0.05)   &      (0.02)   &     (-0.19)   &     (-0.23)   &     (-0.35)   &     (-0.30)   \\
            &               &               &               &               &               &               &               &               \\
Industry FE &           Y   &           Y   &           Y   &           Y   &           Y   &           Y   &           Y   &           Y   \\
Firms       &         426   &         426   &         426   &         426   &         426   &         426   &         426   &         426   \\
R\textsuperscript{2}&       0.140   &       0.139   &       0.163   &       0.162   &       0.181   &       0.177   &       0.163   &       0.164   \\

    \bottomrule
    \end{tabular}


\end{table}
\end{landscape}
\clearpage


\begin{landscape}
\pagestyle{empty}
\begin{table}[htbp]
\caption{\textbf{Credit Spread Response to the Announcement of Moody's Rule Change}
} \label{table_spreadun}
\begin{singlespace}
\begin{small}

\noindent This table contains an event study testing whether treated firms' credit spreads increased around the rule change compared to other firms rated speculative-grade by Moody's. The dependent variable is the change in the firm's credit spread over the period in basis points. T-statistics are shown below the parameter estimates in parentheses and are calculated using robust standard errors. *, **, and *** indicate statistical significance at the $10\%$, $5\%$, and $1\%$ levels, respectively.

\end{small}
\end{singlespace}
\vspace{3mm}
\small
  \centering

\begin{tabular}{lcccccccc}
\\
    \toprule
        &
        \multicolumn{2}{c}{\textbf{Pre-Period}} &
      \multicolumn{4}{c}{\textbf{Event}} &
        \multicolumn{2}{c}{\textbf{Post-Period}}
        \\
                  \cline{2-9}
        \\
    &
    \multicolumn{2}{c}{$\Delta$ \textbf{CS[-30,-2]}} &
      \multicolumn{2}{c}{$\Delta$ \textbf{CS[-1,1]}} &
        \multicolumn{2}{c}{$\Delta$ \textbf{CS[-1,3]}} &
           \multicolumn{2}{c}{$\Delta$ \textbf{CS[4,30]}}
      \\
          \cline{2-9}
            &\multicolumn{1}{c}{(1)}   &\multicolumn{1}{c}{(2)}   &\multicolumn{1}{c}{(3)}   &\multicolumn{1}{c}{(4)}   &\multicolumn{1}{c}{(5)}   &\multicolumn{1}{c}{(6)}   &\multicolumn{1}{c}{(7)}   &\multicolumn{1}{c}{(8)}   \\
\hline
Preferred Dummy&       -16.5   &               &        12.1   &               &        17.7*  &               &        -2.6   &               \\
            &     (-1.22)   &               &      (1.29)   &               &      (1.85)   &               &     (-0.25)   &               \\
Preferred/Capital&               &       -53.1   &               &        32.8   &               &        67.5*  &               &       -15.3   \\
            &               &     (-0.51)   &               &      (1.09)   &               &      (1.75)   &               &     (-0.21)   \\
Profitability&       103.3   &       131.1   &      -225.6*  &      -247.1** &      -191.9   &      -220.6   &      -282.5   &      -279.0   \\
            &      (0.28)   &      (0.36)   &     (-1.95)   &     (-2.27)   &     (-1.26)   &     (-1.52)   &     (-0.91)   &     (-0.92)   \\
Tangibility &       -31.2   &       -29.0   &       -16.6*  &       -18.3*  &        -8.8   &       -10.9   &        -0.6   &        -0.4   \\
            &     (-1.32)   &     (-1.24)   &     (-1.79)   &     (-1.87)   &     (-0.88)   &     (-1.05)   &     (-0.03)   &     (-0.02)   \\
Sales       &        -9.6** &        -9.5** &         3.5** &         3.4***&         4.2***&         4.3***&         3.3   &         3.2   \\
            &     (-2.36)   &     (-2.36)   &      (2.56)   &      (2.62)   &      (2.68)   &      (2.84)   &      (0.75)   &      (0.73)   \\
Market-to-Book&        -4.3   &        -3.9   &         7.4** &         7.0** &         8.3** &         8.0** &        12.4*  &        12.4*  \\
            &     (-0.56)   &     (-0.52)   &      (2.06)   &      (2.01)   &      (2.19)   &      (2.15)   &      (1.93)   &      (1.94)   \\
            &               &               &               &               &               &               &               &               \\
Industry FE &           Y   &           Y   &           Y   &           Y   &           Y   &           Y   &           Y   &           Y   \\
Firms       &         216   &         216   &         222   &         222   &         222   &         222   &         221   &         221   \\
R\textsuperscript{2}&       0.179   &       0.174   &       0.264   &       0.247   &       0.328   &       0.304   &       0.114   &       0.114   \\

    \bottomrule
    \end{tabular}
\end{table}
\end{landscape}
\clearpage
\newpage

\begin{landscape}
\pagestyle{empty}

\begin{table}[htbp]
\caption{\textbf{The Effect of Convertible Debt on Response to Rule Change}} \label{table_convert_unmatched}
\begin{singlespace}
\begin{small}

\noindent   This table contains results testing whether treated firms' responses to the rule change are affected by the amount of convertible debt they had in their capital structure at the time of the rule change, compared to other firms rated speculative-grade by Moody's. Preferred Dummy is an indicator variable that equals one if the firm had preferred stock in its capital structure in the last quarter prior to the rule change. Preferred/Capital is the ratio between the amount of preferred stock and the sum of debt and shareholders' equity in the last quarter prior to the rule change. Convert is the ratio of convertible debt to total debt at the time of the rule change. T-statistics are shown below the parameter estimates in parentheses and are calculated using robust standard errors clustered by firm. *, **, and *** indicate statistical significance at the $10\%$, $5\%$, and $1\%$ levels, respectively.
\end{small}
\end{singlespace}
\vspace{3mm}

\small
  \centering
    \begin{tabular}{lcccccccc}
    \toprule
        \textbf{ } &
      \multicolumn{2}{c}{\textbf{Total Debt}} &  \multicolumn{2}{c}{\textbf{Lev\textsubscript{GAAP}}} &
         \multicolumn{2}{c}{\textbf{PPE}} &  \multicolumn{2}{c}{\textbf{Assets}}
      \\
           \cline{2-9}

            &\multicolumn{1}{c}{(1)}   &\multicolumn{1}{c}{(2)}   &\multicolumn{1}{c}{(3)}   &\multicolumn{1}{c}{(4)}   &\multicolumn{1}{c}{(5)}   &\multicolumn{1}{c}{(6)}   &\multicolumn{1}{c}{(7)}   &\multicolumn{1}{c}{(8)}   \\
\hline
Preferred Dummy x Post&       0.286***&               &       0.064***&               &       0.062   &               &       0.088** &               \\
            &      (3.39)   &               &      (3.39)   &               &      (1.48)   &               &      (2.17)   &               \\
Preferred Dummy x Convert x Post&      -1.869***&               &      -0.470** &               &      -0.338   &               &      -0.794***&               \\
            &     (-4.34)   &               &     (-2.21)   &               &     (-0.77)   &               &     (-2.85)   &               \\
Preferred/Capital x Post&               &       1.760***&               &       0.553***&               &       0.399   &               &       0.613** \\
            &               &      (4.28)   &               &      (4.48)   &               &      (1.46)   &               &      (1.98)   \\
Preferred/Capital x Convert x Post&               &     -11.627***&               &      -4.776***&               &      -0.226   &               &      -2.480   \\
            &               &     (-2.93)   &               &     (-3.41)   &               &     (-0.07)   &               &     (-1.07)   \\
Convert x Post&       0.160   &       0.109   &       0.047   &       0.041   &      -0.227   &      -0.241   &      -0.042   &      -0.071   \\
            &      (0.89)   &      (0.61)   &      (1.28)   &      (1.11)   &     (-1.09)   &     (-1.20)   &     (-0.64)   &     (-1.09)   \\
Profitability&      -1.103***&      -1.101***&      -0.330** &      -0.324*  &      -0.354   &      -0.349   &      -0.119   &      -0.108   \\
            &     (-2.69)   &     (-2.67)   &     (-1.98)   &     (-1.95)   &     (-1.09)   &     (-1.08)   &     (-0.43)   &     (-0.40)   \\
Tangibility &       0.306   &       0.322   &       0.301***&       0.301***&       2.767***&       2.772***&      -0.062   &      -0.055   \\
            &      (0.97)   &      (1.01)   &      (2.68)   &      (2.69)   &      (5.70)   &      (5.71)   &     (-0.27)   &     (-0.24)   \\
Sales       &       0.408***&       0.410***&       0.018   &       0.015   &       0.431***&       0.430***&       0.385***&       0.384***\\
            &      (6.41)   &      (6.48)   &      (1.12)   &      (1.06)   &      (7.28)   &      (7.27)   &      (9.85)   &      (9.93)   \\
Market-to-Book&      -0.147** &      -0.148** &      -0.023   &      -0.024   &      -0.104***&      -0.104***&      -0.136***&      -0.136***\\
            &     (-2.52)   &     (-2.55)   &     (-1.06)   &     (-1.09)   &     (-2.95)   &     (-2.97)   &     (-4.10)   &     (-4.09)   \\
\_cons      &       4.510***&       4.502***&       0.386***&       0.400***&       2.980***&       2.983***&       5.746***&       5.751***\\
            &     (10.86)   &     (10.92)   &      (3.31)   &      (3.66)   &      (5.80)   &      (5.81)   &     (21.55)   &     (21.83)   \\
            &               &               &               &               &               &               &               &               \\
Firm FE     &           Y   &           Y   &           Y   &           Y   &           Y   &           Y   &           Y   &           Y   \\
Quarter FE  &           Y   &           Y   &           Y   &           Y   &           Y   &           Y   &           Y   &           Y   \\
Firm Quarters&        5252   &        5252   &        5284   &        5284   &        5284   &        5284   &        5284   &        5284   \\
R\textsuperscript{2}&       0.117   &       0.113   &       0.046   &       0.052   &       0.379   &       0.379   &       0.305   &       0.304   \\



    \bottomrule
    \end{tabular}
\end{table}
\end{landscape}


\begin{landscape}
\pagestyle{empty}

\begin{table}[htbp]
\caption{\textbf{The Effect of Rule Change on Credit Ratings}} \label{table_cr_all}
\begin{singlespace}
\begin{small}

\noindent   This table contains results testing whether treated firms' ratings change after the rule change, compared to other firms rated speculative-grade by Moody's. The dependent variables, Moody's Rating and S\&P Rating, are categorical variables that take values between 1 and 22 that are mapped from Moody's and S\&P letter ratings where 1 is the highest rating and 22 the lowest. Preferred Dummy is an indicator variable that equals one if the firm had preferred stock in its capital structure in the last quarter prior to the rule change. Preferred/Capital is the ratio between the amount of preferred stock and the sum of debt and shareholders' equity in the last quarter prior to the rule change.  T-statistics are shown below the parameter estimates in parentheses and are calculated using robust standard errors clustered by firm. *, **, and *** indicate statistical significance at the $10\%$, $5\%$, and $1\%$ levels, respectively.
\end{small}
\end{singlespace}
\vspace{3mm}

\small
  \centering
    \begin{tabular}{lcccc}
    \toprule
    \textbf{ } &
      \multicolumn{2}{c}{\textbf{Moody's Rating}} &  \multicolumn{2}{c}{\textbf{S\&P Rating}}
      \\
           \cline{2-5}


            &\multicolumn{1}{c}{(1)}   &\multicolumn{1}{c}{(2)}   &\multicolumn{1}{c}{(3)}   &\multicolumn{1}{c}{(4)}   \\
\hline
Preferred Dummy x Post&       0.003   &               &      -0.065   &               \\
            &      (0.02)   &               &     (-0.46)   &               \\
Preferred/Capital x Post&               &      -1.069   &               &      -0.299   \\
            &               &     (-1.21)   &               &     (-0.38)   \\
Profitability&      -1.026   &      -1.047   &      -1.794   &      -1.794   \\
            &     (-1.17)   &     (-1.20)   &     (-1.56)   &     (-1.56)   \\
Tangibility &       0.214   &       0.193   &       0.274   &       0.271   \\
            &      (0.35)   &      (0.32)   &      (0.58)   &      (0.57)   \\
Sales       &      -0.188** &      -0.184** &      -0.428***&      -0.428***\\
            &     (-2.16)   &     (-2.11)   &     (-4.35)   &     (-4.35)   \\
Market-to-Book&      -0.315***&      -0.315***&      -0.399***&      -0.398***\\
            &     (-3.00)   &     (-2.99)   &     (-3.99)   &     (-3.98)   \\
\_cons      &      15.604***&      15.594***&      16.296***&      16.297***\\
            &     (24.43)   &     (24.42)   &     (25.77)   &     (25.77)   \\
            &               &               &               &               \\
Firm FE     &           Y   &           Y   &           Y   &           Y   \\
Quarter FE  &           Y   &           Y   &           Y   &           Y   \\
Firm Quarters&        5051   &        5051   &        4835   &        4835   \\
R\textsuperscript{2}&       0.026   &       0.027   &       0.085   &       0.085   \\



    \bottomrule
    \end{tabular}
\end{table}
\end{landscape}

\clearpage
\newpage


\subsection*{Additional Robustness Tests}

\begin{table}[htbp]
\caption{\textbf{The Effect of Rule Change on GAAP Leverage Using IG Firms as a Control Group}} \label{table_diff_c}
\begin{singlespace}
\begin{small}

\noindent This table contains results testing whether speculative-rated firms with preferred stock increase their GAAP leverage after the rule change compared to investment-grade firms with preferred stock. The dependent variable is $Lev_{GAAP}$. Preferred Dummy is an indicator variable that equals one if the firm had preferred stock in its capital structure in the last quarter prior to the rule change. $Junk$ is an indicator variable that equals one if the firm was rated speculative-grade by Moody's at the time of the rule change. In Columns 1 - 2, the sample is restricted to firms with preferred stock at the time of the rule change, and in Columns 3 - 4, the sample includes all treated and untreated firms. T-statistics are shown below the parameter estimates in parentheses and are calculated using robust standard errors clustered by firm. *, **, and *** indicate statistical significance at the $10\%$, $5\%$, and $1\%$ levels, respectively.

\end{small}
\end{singlespace}
\vspace{3mm}
\small
  \centering
    \begin{tabular}{lcccc}
    \toprule
    \textbf{ } &
      \multicolumn{4}{c}{\textbf{Lev\textsubscript{GAAP}}}
      \\
      \cline{2-5}

                  &\multicolumn{1}{c}{(1)}   &\multicolumn{1}{c}{(2)}   &\multicolumn{1}{c}{(3)}   &\multicolumn{1}{c}{(4)}   \\
\hline
Junk x Post &       0.109***&       0.058***&       0.050***&      -0.003   \\
            &      (4.84)   &      (2.62)   &      (7.41)   &     (-0.51)   \\
Preferred x Post x Junk&               &               &       0.060** &       0.059***\\
            &               &               &      (2.56)   &      (2.58)   \\
Preferred x Post&               &               &      -0.007   &      -0.012   \\
            &               &               &     (-0.62)   &     (-1.09)   \\
Profitability&               &       0.043***&               &       0.036***\\
            &               &      (3.36)   &               &      (7.49)   \\
Tangibility &               &      -0.005   &               &       0.101***\\
            &               &     (-0.05)   &               &      (3.00)   \\
Sales       &               &       0.002   &               &       0.020***\\
            &               &      (0.21)   &               &      (3.61)   \\
Market-to-Book&               &      -0.000   &               &      -0.000***\\
            &               &     (-0.11)   &               &     (-3.00)   \\
Constant    &       0.356***&       0.331***&       0.359***&       0.212***\\
            &    (265.40)   &      (7.12)   &    (400.94)   &      (9.28)   \\
            &               &               &               &               \\
Sample      &   Preferred   &   Preferred   &         All   &         All   \\
Firm FE     &           Y   &           Y   &           Y   &           Y   \\
Quarter FE  &           Y   &           Y   &           Y   &           Y   \\
Quarter x Cohort FE&           N   &           N   &           N   &           N   \\
Firm Quarters&        5761   &        4922   &       55732   &       46543   \\
R\textsuperscript{2}&       0.005   &       0.016   &       0.002   &       0.020   \\



    \bottomrule
    \end{tabular}
\end{table}
\clearpage

\newpage

\begin{table}[htbp]
\caption{\textbf{The Effect of Rule Change on GAAP Leverage: Debt/Assets}} \label{table_bla_c}
\begin{singlespace}
\begin{small}

\noindent This table contains results testing whether treated firms increase their GAAP leverage after the rule. The dependent variables are Debt/Assets\textsubscript{GAAP} and Debt/Assets\textsubscript{Moody's}. Preferred Dummy is an indicator variable that equals one if the firm had preferred stock in its capital structure in the last quarter prior to the rule change. Preferred/Assets is preferred stock over assets at the time of the rule change. The samples include matched and unmatched firms. T-statistics are shown below the parameter estimates in parentheses and are calculated using robust standard errors clustered by firm. *, **, and *** indicate statistical significance at the $10\%$, $5\%$, and $1\%$ levels, respectively.

\end{small}
\end{singlespace}
\vspace{3mm}
\small
  \centering
    \begin{tabular}{lcccc}
    \toprule
    \textbf{ } &
      \multicolumn{2}{c}{\textbf{Debt/Assets\textsubscript{GAAP}}} &
       \multicolumn{2}{c}{\textbf{Debt/Assets\textsubscript{Moody's}}}
      \\
      \cline{2-5}

              &\multicolumn{1}{c}{(1)}   &\multicolumn{1}{c}{(2)}   &\multicolumn{1}{c}{(3)}   &\multicolumn{1}{c}{(4)}   \\
\hline
Preferred Dummy x Post&       0.036** &               &       0.003   &               \\
            &      (2.43)   &               &      (0.17)   &               \\
Preferred/Assets x Post&               &       0.368*  &               &      -0.179   \\
            &               &      (1.93)   &               &     (-0.92)   \\
Profitability&      -0.283***&      -0.287***&      -0.294***&      -0.291***\\
            &     (-5.09)   &     (-5.17)   &     (-5.30)   &     (-5.27)   \\
Tangibility &       0.289***&       0.292***&       0.295***&       0.290***\\
            &      (2.69)   &      (2.71)   &      (2.77)   &      (2.70)   \\
Sales       &       0.001   &       0.002   &       0.000   &       0.001   \\
            &      (0.07)   &      (0.13)   &      (0.00)   &      (0.06)   \\
Market-to-Book&       0.028   &       0.026   &       0.022   &       0.025   \\
            &      (0.92)   &      (0.84)   &      (0.70)   &      (0.78)   \\
\_cons      &       0.259** &       0.255** &       0.276** &       0.272** \\
            &      (2.27)   &      (2.25)   &      (2.44)   &      (2.40)   \\
            &               &               &               &               \\
Firm FE     &           Y   &           Y   &           Y   &           Y   \\
Quarter x Cohort FE&           Y   &           Y   &           Y   &           Y   \\
Firm Quarters&        2454   &        2454   &        2454   &        2454   \\
R\textsuperscript{2}&       0.078   &       0.074   &       0.065   &       0.068   \\



    \bottomrule
    \end{tabular}
\end{table}
\clearpage

\newpage

\begin{table}[htbp]
\caption{\textbf{Market Trends around the Announcement of the Rule Change}} \label{table_carnews}
\begin{singlespace}
\begin{small}

\noindent This table shows the cumulative return of the S\&P 500 and the 10-year treasury cumulative yield change in the trading days [-1;+1],[-3;+3], and [-5;+5] centered around Moody's announcement of the rule change (July 31, 2013).

\end{small}
\end{singlespace}
\vspace{3mm}
\small
  \centering
    \begin{tabular}{lccc}
    \toprule
    \textbf{ } &
      \multicolumn{3}{c}{\textbf{Days Around Moody's}}\\
          \textbf{ } &
      \multicolumn{3}{c}{\textbf{Rule Change Announcement}}
      \\
      \cline{2-4}

            &\multicolumn{1}{c}{[-1,+1]}   &\multicolumn{1}{c}{[-3,+3]}   &\multicolumn{1}{c}{[-5,+5]}   \\
\hline
S\&P 500 Cumulative Return (\%)&       1.24   &       0.92   &       0.29   \\
10 Year Treasury Cumulative Yield Change (bps)&       12.00   &       7.90   &       1.20   \\


    \bottomrule
    \end{tabular}
\end{table}
\clearpage

\newpage


\begin{table}[htbp]
\caption{\textbf{Determinants of Preferred Stock}} \label{table_prefdet}
\begin{singlespace}
\begin{small}

\noindent This table contains results testing the determinants of firms using preferred stock at the time of the rule change. T-statistics are shown below the parameter estimates in parentheses and are calculated using robust standard errors clustered by firm. *, **, and *** indicate statistical significance at the $10\%$, $5\%$, and $1\%$ levels, respectively.

\end{small}
\end{singlespace}
\vspace{3mm}
\small
  \centering
    \begin{tabular}{lccccc}
    \toprule
    \textbf{ } &
      \multicolumn{5}{c}{\textbf{Preferred Dummy}}
      \\
      \cline{2-6}

              &\multicolumn{1}{c}{(1)}   &\multicolumn{1}{c}{(2)}   &\multicolumn{1}{c}{(3)}   &\multicolumn{1}{c}{(4)}   &\multicolumn{1}{c}{(5)}   \\
\hline
Profitability&      -0.025*  &      -0.709   &       0.144   &      -0.016   &       0.457   \\
            &     (-1.92)   &     (-0.86)   &      (0.11)   &     (-1.15)   &      (0.18)   \\
Tangibility &       0.096** &      -0.126   &      -0.067   &       0.098** &      -0.342   \\
            &      (2.38)   &     (-0.83)   &     (-0.32)   &      (2.17)   &     (-0.92)   \\
Sales       &      -0.009   &      -0.022   &      -0.018   &      -0.012*  &      -0.023   \\
            &     (-1.53)   &     (-0.90)   &     (-0.46)   &     (-1.83)   &     (-0.33)   \\
Market-to-Book&      -0.001***&      -0.024   &      -0.017   &      -0.001***&      -0.241   \\
            &     (-4.36)   &     (-0.91)   &     (-0.70)   &     (-4.60)   &     (-1.53)   \\
Leverage    &       0.064***&       0.025   &      -0.034   &       0.060***&       0.069   \\
            &      (3.77)   &      (0.41)   &     (-0.37)   &      (3.10)   &      (0.38)   \\
Market Cap  &       0.000   &      -0.003   &      -0.001   &       0.000   &      -0.000   \\
            &      (1.59)   &     (-0.82)   &     (-1.14)   &      (0.43)   &     (-0.06)   \\
Log(PPE)    &      -0.012*  &       0.025   &       0.016   &      -0.011   &       0.049   \\
            &     (-1.94)   &      (0.80)   &      (0.33)   &     (-1.57)   &      (0.66)   \\
Log(Assets) &       0.004   &       0.017   &       0.043   &      -0.003   &      -0.025   \\
            &      (0.57)   &      (0.51)   &      (0.79)   &     (-0.31)   &     (-0.26)   \\
            &               &               &               &               &               \\
Sample      &         All   &        Junk   &          IG   &     Unrated   &     Matched   \\
Industry FE &           Y   &           Y   &           Y   &           Y   &           Y   \\
Observations&        3791   &         539   &         391   &        2857   &         208   \\
R\textsuperscript{2}&       0.045   &       0.104   &       0.161   &       0.055   &       0.031   \\



    \bottomrule
    \end{tabular}
\end{table}
\clearpage

\newpage


\begin{table}[htbp]
\caption{\textbf{Determinants of Convertibles}} \label{table_convdet}
\begin{singlespace}
\begin{small}

\noindent This table contains results testing the determinants of firms using convertible debt at the time of the rule change. The dependent variable is the ratio of convertible debt to total debt. T-statistics are shown below the parameter estimates in parentheses and are calculated using robust standard errors clustered by firm. *, **, and *** indicate statistical significance at the $10\%$, $5\%$, and $1\%$ levels, respectively.

\end{small}
\end{singlespace}
\vspace{3mm}
\small
  \centering
    \begin{tabular}{lccccc}
    \toprule
    \textbf{ } &
      \multicolumn{5}{c}{\textbf{Convertible}}
      \\
      \cline{2-6}

              &\multicolumn{1}{c}{(1)}   &\multicolumn{1}{c}{(2)}   &\multicolumn{1}{c}{(3)}   &\multicolumn{1}{c}{(4)}   &\multicolumn{1}{c}{(5)}   \\
\hline
Profitability&      -0.244***&      -0.257*  &       0.012   &      -0.243***&       0.078   \\
            &     (-4.65)   &     (-1.91)   &      (0.14)   &     (-3.93)   &      (0.32)   \\
Tangibility &      -0.177   &       0.004   &       0.009   &      -0.254   &      -0.026   \\
            &     (-1.13)   &      (0.16)   &      (0.70)   &     (-1.26)   &     (-0.77)   \\
Sales       &      -0.007   &       0.002   &      -0.002   &      -0.012   &      -0.011   \\
            &     (-0.28)   &      (0.45)   &     (-0.75)   &     (-0.41)   &     (-1.58)   \\
Market-to-Book&      -0.003***&       0.005   &       0.002   &      -0.003***&       0.013   \\
            &     (-3.43)   &      (1.09)   &      (1.18)   &     (-2.97)   &      (0.85)   \\
Leverage    &       0.241***&       0.023** &       0.001   &       0.297***&       0.034** \\
            &      (3.68)   &      (2.37)   &      (0.13)   &      (3.45)   &      (2.05)   \\
Market Cap  &       0.001   &      -0.001   &      -0.000   &       0.002   &      -0.001   \\
            &      (1.01)   &     (-0.94)   &     (-0.84)   &      (0.51)   &     (-1.45)   \\
Log(PPE)    &       0.013   &      -0.002   &      -0.007** &       0.020   &      -0.001   \\
            &      (0.51)   &     (-0.42)   &     (-2.15)   &      (0.66)   &     (-0.10)   \\
Log(Assets) &      -0.036   &       0.003   &       0.009** &      -0.044   &       0.016*  \\
            &     (-1.20)   &      (0.60)   &      (2.47)   &     (-1.17)   &      (1.88)   \\
            &               &               &               &               &               \\
Sample      &         All   &        Junk   &          IG   &     Unrated   &     Matched   \\
Industry FE &           Y   &           Y   &           Y   &           Y   &           Y   \\
Observations&        3713   &         535   &         388   &        2785   &         204   \\
R\textsuperscript{2}&       0.026   &       0.107   &       0.194   &       0.030   &       0.140   \\



    \bottomrule
    \end{tabular}
\end{table}
\clearpage

\newpage


\begin{table}[htbp]
\caption{\textbf{The Effect of Convertible Debt on Response to Rule Change Using IG Firms as Control Group }} \label{table_convertIG}
\begin{singlespace}
\begin{small}

\noindent This table contains results testing whether the response of speculative-rated firms with preferred stock to the rule change is affected by the amount of convertible debt they had in their capital structure at the time of the rule change, relative to investment-grade firms with preferred stock. Junk is a dummy equal to one if the firm is rated speculative-grade by Moody's, and 0 if it is rated investment-grade. Convert is the ratio of convertible debt to total debt at the time of the rule change. T-statistics are shown below the parameter estimates in parentheses and are calculated using robust standard errors clustered by firm. *, **, and *** indicate statistical significance at the 10


\end{small}
\end{singlespace}
\vspace{3mm}
\small
  \centering
    \begin{tabular}{lcccc}
    \toprule
    \textbf{ } &
      \multicolumn{1}{c}{\textbf{Total Debt}} &  \multicolumn{1}{c}{\textbf{Lev\textsubscript{GAAP}}} &
         \multicolumn{1}{c}{\textbf{PPE}} &  \multicolumn{1}{c}{\textbf{Assets}}
      \\
      \cline{2-5}

              &\multicolumn{1}{c}{(1)}   &\multicolumn{1}{c}{(2)}   &\multicolumn{1}{c}{(3)}   &\multicolumn{1}{c}{(4)}   \\
\hline
Convert x Post x Junk&      -1.426***&      -0.381*  &      -0.871*  &      -0.935***\\
            &     (-3.28)   &     (-1.84)   &     (-1.93)   &     (-3.08)   \\
Convert x Post&      -0.183   &      -0.038   &       0.489** &       0.232*  \\
            &     (-0.91)   &     (-0.95)   &      (1.97)   &      (1.94)   \\
Post x Junk &       0.144   &       0.073***&       0.109*  &       0.046   \\
            &      (1.42)   &      (3.17)   &      (1.90)   &      (0.93)   \\
Profitability&       0.055   &       0.043***&       0.321***&       0.439***\\
            &      (1.20)   &      (3.32)   &      (3.74)   &      (8.11)   \\
Tangibility &       0.754** &       0.000   &       4.830***&       0.244   \\
            &      (2.24)   &      (0.00)   &      (9.39)   &      (0.80)   \\
Sales       &       0.326***&       0.004   &       0.388***&       0.403***\\
            &      (4.09)   &      (0.33)   &      (4.94)   &      (5.30)   \\
Market-to-Book&      -0.000   &      -0.000   &      -0.005***&      -0.008***\\
            &     (-0.37)   &     (-0.12)   &     (-3.86)   &     (-7.09)   \\
            &               &               &               &               \\
Firm FE     &           Y   &           Y   &           Y   &           Y   \\
Quarter FE  &           Y   &           Y   &           Y   &           Y   \\
Firm Quarters&        4194   &        4885   &        4574   &        4885   \\
R\textsuperscript{2}&       0.038   &       0.017   &       0.416   &       0.519   \\


    \bottomrule
    \end{tabular}
\end{table}
\clearpage


\begin{table}[htbp]
\caption{\textbf{The Effect of Rule Change on Debt Levels - One Matched Firm with No Replacement}} \label{table_debt_alt}
\begin{singlespace}
\begin{small}

\noindent This table contains results testing whether treated firms increase their debt levels after the rule change compared to matched firms. Each treated firm is matched without replacement to a single untreated firm with the most similar characteristics, as detailed in section \ref{sec_experiment}. Coefficients are estimated using Equation (\ref{eq:single}). The dependent variables are total debt (Columns 1 and 2), long-term debt (Columns 3 and 4) and short-term debt (Columns 5 and 6) and are all in logs. Preferred Dummy is an indicator variable that equals one if the firm had preferred stock in its capital structure in the last quarter prior to the rule change. Preferred/Capital is the ratio between the amount of preferred stock and the sum of debt and shareholders' equity in the last quarter prior to the rule change. T-statistics are shown below the parameter estimates in parentheses and are calculated using robust standard errors clustered by firm. *, **, and *** indicate statistical significance at the $10\%$, $5\%$, and $1\%$ levels, respectively.

\end{small}
\end{singlespace}
\vspace{3mm}
\small
  \centering
    \begin{tabular}{lcccccc}
    \toprule
    \textbf{} &
      \multicolumn{2}{c}{\textbf{Total Debt}} &  \multicolumn{2}{c}{\textbf{LT Debt}} &
      \multicolumn{2}{c}{\textbf{ST Debt}}
      \\
      \cline{2-7}
            &\multicolumn{1}{c}{(1)}   &\multicolumn{1}{c}{(2)}   &\multicolumn{1}{c}{(3)}   &\multicolumn{1}{c}{(4)}   &\multicolumn{1}{c}{(5)}   &\multicolumn{1}{c}{(6)}   \\
\hline
Preferred Dummy x Post&       0.307***&               &       0.282***&               &       0.432   &               \\
            &      (3.83)   &               &      (3.75)   &               &      (1.17)   &               \\
Preferred/Capital x Post&               &       1.160** &               &       1.418** &               &       4.103   \\
            &               &      (2.24)   &               &      (2.52)   &               &      (0.93)   \\
Profitability&       0.035   &      -0.063   &      -1.231** &      -1.426** &       0.718   &       0.988   \\
            &      (0.07)   &     (-0.11)   &     (-2.14)   &     (-2.31)   &      (0.25)   &      (0.36)   \\
Tangibility &      -0.591   &      -0.623   &      -1.008   &      -1.089   &      -3.114   &      -3.369   \\
            &     (-0.98)   &     (-0.91)   &     (-1.57)   &     (-1.52)   &     (-1.20)   &     (-1.30)   \\
Sales       &       0.477***&       0.496***&       0.462***&       0.472***&      -0.272   &      -0.294   \\
            &      (3.74)   &      (3.92)   &      (3.73)   &      (3.86)   &     (-0.60)   &     (-0.65)   \\
Market-to-Book&       0.433   &       0.472   &       0.346   &       0.352   &      -0.775*  &      -0.743*  \\
            &      (1.35)   &      (1.31)   &      (0.93)   &      (0.85)   &     (-1.89)   &     (-1.75)   \\
\_cons      &       3.878***&       3.810***&       4.209***&       4.226***&       6.606** &       6.822** \\
            &      (4.22)   &      (3.97)   &      (4.86)   &      (4.76)   &      (2.03)   &      (2.07)   \\
            &               &               &               &               &               &               \\
Firm FE     &           Y   &           Y   &           Y   &           Y   &           Y   &           Y   \\
Quarter x Cohort FE&           Y   &           Y   &           Y   &           Y   &           Y   &           Y   \\
Firm Quarters&         984   &         984   &         958   &         958   &         600   &         600   \\
R\textsuperscript{2}&       0.269   &       0.206   &       0.237   &       0.193   &       0.046   &       0.051   \\



    \bottomrule
    \end{tabular}
\end{table}
\clearpage


\begin{table}[htbp]
\caption{\textbf{The Effect of Rule Change on Leverage - One Matched Firm with No Replacement}
} \label{table_lev_alt}
\begin{singlespace}
\begin{small}

\noindent This table contains results testing whether treated firms increase their GAAP leverage after the rule change compared to matched firms. Each treated firm is matched without replacement to a single untreated firm with the most similar characteristics, as detailed in section \ref{sec_experiment}. Coefficients are estimated using Equation (\ref{eq:single}). The dependent variables are Lev\textsubscript{GAAP} and $Lev_{Moody's}$. Preferred Dummy is an indicator variable that equals one if the firm had preferred stock in its capital structure in the last quarter prior to the rule change. Preferred/Capital is the ratio between the amount of preferred stock and the sum of debt and shareholders' equity in the last quarter prior to the rule change. T-statistics are shown below the parameter estimates in parentheses and are calculated using robust standard errors clustered by firm. *, **, and *** indicate statistical significance at the $10\%$, $5\%$, and $1\%$ levels, respectively.

\end{small}
\end{singlespace}
\vspace{3mm}
\small
  \centering
   \scalebox{1}{ \begin{tabular}{lcccc}
    \toprule
    \textbf{ } &
      \multicolumn{2}{c}{\textbf{Lev\textsubscript{GAAP}}}  &
        \multicolumn{2}{c}{\textbf{Lev\textsubscript{Moody's}}}
      \\
          \cline{2-5}

              &\multicolumn{1}{c}{(1)}   &\multicolumn{1}{c}{(2)}   &\multicolumn{1}{c}{(3)}   &\multicolumn{1}{c}{(4)}   \\
\hline
Preferred Dummy x Post&       0.046** &               &      -0.000   &               \\
            &      (2.31)   &               &     (-0.02)   &               \\
Preferred/Capital x Post&               &       0.261   &               &      -0.277   \\
            &               &      (1.24)   &               &     (-1.18)   \\
Profitability&      -0.076   &      -0.095   &      -0.126   &      -0.114   \\
            &     (-0.20)   &     (-0.25)   &     (-0.32)   &     (-0.30)   \\
Tangibility &      -0.246   &      -0.247   &      -0.270   &      -0.275   \\
            &     (-1.34)   &     (-1.36)   &     (-1.46)   &     (-1.50)   \\
Sales       &      -0.012   &      -0.010   &      -0.014   &      -0.013   \\
            &     (-0.44)   &     (-0.34)   &     (-0.51)   &     (-0.47)   \\
Market-to-Book&      -0.043   &      -0.040   &      -0.051   &      -0.043   \\
            &     (-1.11)   &     (-0.93)   &     (-1.08)   &     (-0.98)   \\
\_cons      &       0.799***&       0.788***&       0.853***&       0.850***\\
            &      (4.34)   &      (4.27)   &      (4.66)   &      (4.61)   \\
            &               &               &               &               \\
Firm FE     &           Y   &           Y   &           Y   &           Y   \\
Quarter x Cohort FE&           Y   &           Y   &           Y   &           Y   \\
Firm Quarters&         996   &         996   &         996   &         996   \\
R\textsuperscript{2}&       0.062   &       0.041   &       0.033   &       0.056   \\



    \bottomrule
    \end{tabular}}
\end{table}


\clearpage

\newpage


\begin{table}[htbp]
\caption{\textbf{The Effect of Rule Change on the Balance Sheet - One Matched Firm with No Replacement}} \label{table_bs_alt}
\begin{singlespace}
\begin{small}

\noindent This table contains results testing whether treated firms increase the size of their balance sheets after the rule change compared to matched firms. Each treated firm is matched without replacement to a single untreated firm with the most similar characteristics, as detailed in section \ref{sec_experiment}.  Coefficients are estimated using Equation (\ref{eq:single}). The dependent variables are Capex, PPE and Assets, all in logs. Preferred Dummy is an indicator variable that equals one if the firm had preferred stock in its capital structure in the last quarter prior to the rule change. Preferred/Capital is the ratio between the amount of preferred stock and the sum of debt and shareholders' equity in the last quarter prior to the rule change. T-statistics are shown below the parameter estimates in parentheses and are calculated using robust standard errors clustered by firm. *, **, and *** indicate statistical significance at the $10\%$, $5\%$, and $1\%$ levels, respectively.

\end{small}
\end{singlespace}
\vspace{3mm}
\small
  \centering
    \begin{tabular}{lcccccc}
    \toprule
    \textbf{ } &
      \multicolumn{2}{c}{\textbf{Capex}} &  \multicolumn{2}{c}{\textbf{PPE}} &
      \multicolumn{2}{c}{\textbf{Assets}}
      \\
      \cline{2-7}
              &\multicolumn{1}{c}{(1)}   &\multicolumn{1}{c}{(2)}   &\multicolumn{1}{c}{(3)}   &\multicolumn{1}{c}{(4)}   &\multicolumn{1}{c}{(5)}   &\multicolumn{1}{c}{(6)}   \\
\hline
Preferred Dummy x Post&       0.157*  &               &       0.085** &               &       0.099***&               \\
            &      (1.91)   &               &      (2.59)   &               &      (2.64)   &               \\
Preferred/Capital x Post&               &       0.995*  &               &       0.345   &               &       0.527*  \\
            &               &      (1.89)   &               &      (1.47)   &               &      (1.82)   \\
Profitability&       1.061*  &       0.993*  &       0.288   &       0.259   &       0.287   &       0.248   \\
            &      (1.96)   &      (1.73)   &      (1.19)   &      (1.07)   &      (1.22)   &      (1.02)   \\
Tangibility &       2.120   &       2.121   &       1.988***&       1.983***&      -0.049   &      -0.053   \\
            &      (1.44)   &      (1.41)   &      (5.47)   &      (5.04)   &     (-0.12)   &     (-0.12)   \\
Sales       &       0.593***&       0.601***&       0.370***&       0.376***&       0.395***&       0.401***\\
            &      (3.63)   &      (3.71)   &      (4.64)   &      (4.63)   &      (4.69)   &      (4.71)   \\
Market-to-Book&       0.494***&       0.502***&      -0.049   &      -0.039   &      -0.055   &      -0.046   \\
            &      (2.94)   &      (2.65)   &     (-0.47)   &     (-0.36)   &     (-0.48)   &     (-0.40)   \\
\_cons      &      -1.456   &      -1.491   &       3.623***&       3.602***&       5.606***&       5.583***\\
            &     (-1.24)   &     (-1.26)   &      (6.56)   &      (6.36)   &      (9.99)   &      (9.75)   \\
            &               &               &               &               &               &               \\
Firm FE     &           Y   &           Y   &           Y   &           Y   &           Y   &           Y   \\
Quarter x Cohort FE&           Y   &           Y   &           Y   &           Y   &           Y   &           Y   \\
Firm Quarters&         994   &         994   &         996   &         996   &         996   &         996   \\
R\textsuperscript{2}&       0.153   &       0.148   &       0.358   &       0.337   &       0.339   &       0.318   \\



    \bottomrule
    \end{tabular}
\end{table}
\clearpage

\newpage


\begin{table}[htbp]
\caption{\textbf{The Effect of Rule Change on the PP\&E - Inverse Hyperbolic Sine Transformation}} \label{table_bs_sine}
\begin{singlespace}
\begin{small}

\noindent This table contains results testing whether treated firms increase the size of their balance sheets after the rule change compared to matched firms. The dependent variable is the IHS transformation of PP\&E. Preferred Dummy is an indicator variable that equals one if the firm had preferred stock in its capital structure in the last quarter prior to the rule change. Preferred/Capital is the ratio between the amount of preferred stock and the sum of debt and shareholders' equity in the last quarter prior to the rule change. T-statistics are shown below the parameter estimates in parentheses and are calculated using robust standard errors clustered by firm. *, **, and *** indicate statistical significance at the $10\%$, $5\%$, and $1\%$ levels, respectively.

\end{small}
\end{singlespace}
\vspace{3mm}
\small
  \centering
    \begin{tabular}{lcc}
    \toprule
    \textbf{ } &
      \multicolumn{2}{c}{\textbf{IHS PPE}}
      \\
      \cline{2-3}
              &\multicolumn{1}{c}{(1)}   &\multicolumn{1}{c}{(2)}   \\
\hline
Preferred Dummy x Post&       0.081** &               \\
            &      (2.31)   &               \\
Preferred/Capital x Post&               &       0.473*  \\
            &               &      (1.74)   \\
Profitability&       0.332** &       0.327** \\
            &      (2.47)   &      (2.39)   \\
Tangibility &       1.966***&       1.967***\\
            &      (4.67)   &      (4.63)   \\
Sales       &       0.352***&       0.354***\\
            &      (7.38)   &      (7.48)   \\
Market-to-Book&      -0.144** &      -0.146** \\
            &     (-2.21)   &     (-2.22)   \\
\_cons      &       4.579***&       4.573***\\
            &     (15.78)   &     (15.67)   \\
            &               &               \\
Firm FE     &           Y   &           Y   \\
Quarter x Cohort FE&           Y   &           Y   \\
Firm Quarters&        2454   &        2454   \\
R\textsuperscript{2}&       0.337   &       0.334   \\



    \bottomrule
    \end{tabular}
\end{table}
\clearpage

\newpage

\clearpage
\newpage

\begin{landscape}
\pagestyle{empty}

\begin{table}[htbp]
\caption{\textbf{The Effect of Convertible Debt on Response to Rule Change - Convertible Dummy}} \label{table_convert_dummy}
\begin{singlespace}
\begin{small}

\noindent   This table contains results testing whether treated firms' responses to the rule change are affected by the amount of convertible debt they had in their capital structure at the time of the rule change. Preferred Dummy is an indicator variable that equals one if the firm had preferred stock in its capital structure in the last quarter prior to the rule change. Preferred/Capital is the ratio between the amount of preferred stock and the sum of debt and shareholders' equity in the last quarter prior to the rule change. Convert Dummy is a dummy variable equal 1 if the firm had any convertible debt at the time of the rule change. T-statistics are shown below the parameter estimates in parentheses and are calculated using robust standard errors clustered by firm. *, **, and *** indicate statistical significance at the $10\%$, $5\%$, and $1\%$ levels, respectively.
\end{small}
\end{singlespace}
\vspace{3mm}

\small
  \centering
    \begin{tabular}{lcccccccc}
    \toprule
        \textbf{ } &
      \multicolumn{2}{c}{\textbf{Total Debt}} &  \multicolumn{2}{c}{\textbf{Lev\textsubscript{GAAP}}} &
         \multicolumn{2}{c}{\textbf{PPE}} &  \multicolumn{2}{c}{\textbf{Assets}}
      \\
           \cline{2-9}

            &\multicolumn{1}{c}{(1)}   &\multicolumn{1}{c}{(2)}   &\multicolumn{1}{c}{(3)}   &\multicolumn{1}{c}{(4)}   &\multicolumn{1}{c}{(5)}   &\multicolumn{1}{c}{(6)}   &\multicolumn{1}{c}{(7)}   &\multicolumn{1}{c}{(8)}   \\
\hline
Preferred Dummy x Post&       0.306***&               &       0.056***&               &       0.129***&               &       0.117** &               \\
            &      (3.79)   &               &      (2.79)   &               &      (3.07)   &               &      (2.48)   &               \\
Preferred Dummy x Convert x Post&      -0.290** &               &      -0.087*  &               &      -0.177** &               &      -0.120   &               \\
            &     (-2.34)   &               &     (-1.88)   &               &     (-2.22)   &               &     (-1.51)   &               \\
Preferred/Capital x Post&               &       1.527***&               &       0.304** &               &       0.631** &               &       0.619*  \\
            &               &      (3.60)   &               &      (2.01)   &               &      (2.09)   &               &      (1.77)   \\
Preferred/Capital x Convert x Post&               &      -2.377** &               &      -0.759** &               &      -1.358   &               &      -0.485   \\
            &               &     (-2.18)   &               &     (-2.23)   &               &     (-1.55)   &               &     (-0.72)   \\
Convert x Post&       0.011   &       0.002   &      -0.004   &      -0.009   &       0.035   &       0.024   &       0.020   &       0.010   \\
            &      (0.15)   &      (0.02)   &     (-0.14)   &     (-0.37)   &      (0.89)   &      (0.62)   &      (0.50)   &      (0.25)   \\
Profitability&      -0.025   &      -0.034   &      -0.224** &      -0.227** &       0.329** &       0.325** &       0.446***&       0.441***\\
            &     (-0.11)   &     (-0.14)   &     (-2.08)   &     (-2.10)   &      (2.48)   &      (2.38)   &      (3.57)   &      (3.44)   \\
Tangibility &       0.196   &       0.187   &       0.365** &       0.364** &       1.948***&       1.946***&      -0.421   &      -0.415   \\
            &      (0.43)   &      (0.40)   &      (2.38)   &      (2.37)   &      (4.62)   &      (4.54)   &     (-1.27)   &     (-1.23)   \\
Sales       &       0.348***&       0.359***&      -0.014   &      -0.012   &       0.346***&       0.351***&       0.335***&       0.340***\\
            &      (5.61)   &      (5.91)   &     (-0.66)   &     (-0.57)   &      (7.11)   &      (7.38)   &      (6.54)   &      (6.74)   \\
Market-to-Book&      -0.045   &      -0.045   &      -0.072** &      -0.072** &      -0.153** &      -0.151** &      -0.145** &      -0.145** \\
            &     (-0.25)   &     (-0.23)   &     (-2.07)   &     (-2.00)   &     (-2.37)   &     (-2.30)   &     (-2.17)   &     (-2.14)   \\
            &               &               &               &               &               &               &               &               \\
Firm FE     &           Y   &           Y   &           Y   &           Y   &           Y   &           Y   &           Y   &           Y   \\
Quarter x Cohort FE&           Y   &           Y   &           Y   &           Y   &           Y   &           Y   &           Y   &           Y   \\
Firm Quarters&        2448   &        2448   &        2454   &        2454   &        2454   &        2454   &        2454   &        2454   \\
R\textsuperscript{2}&       0.137   &       0.119   &       0.062   &       0.057   &       0.344   &       0.337   &       0.306   &       0.299   \\



    \bottomrule
    \end{tabular}
\end{table}
\end{landscape}


\clearpage


\subsection*{Sample Language in Filings from Treated Firms}\label{filing_lang}
In this section, we include sample language from filings of treated firms relating to rating-based triggers.

\begin{displayquote}
\begin{small}
\begin{singlespace}
``Any further downgrade of Alcoa's credit ratings could limit Alcoa's ability to obtain future financing, increase its borrowing costs, increase the pricing of its credit facilities, adversely affect the market price of its securities, trigger letter of credit or other collateral postings, or otherwise impair its business, financial condition, and results of operations.''
\href{https://www.sec.gov/Archives/edgar/data/4281/000119312513404066/d593859d10q.htm}{Alcoa 10-q}
\end{singlespace}
\end{small}
\end{displayquote}

\begin{displayquote}
\begin{small}
\begin{singlespace}
``Chesapeake has significant flexibility with regard to releases and/or substitutions of pledged reserves, provided that certain requirements are met including maintaining specified collateral coverage ratios as well as maintaining credit ratings with either of the designated rating agencies at or above current levels'' \href{https://www.sec.gov/Archives/edgar/data/895126/000089512613000254/chk-2013063010xq.htm}{Chesapeake Energy 10-q}


\end{singlespace}
\end{small}
\end{displayquote}


\begin{displayquote}
\begin{small}
\begin{singlespace}
``If there is a change of control of the Company and if the Company’s corporate credit rating is withdrawn or downgraded to a certain level (together constituting a “change of control event”), the dividends on the Series A Preferred Shares will increase to 10.75\% per year.'' \href{https://www.sec.gov/Archives/edgar/data/799292/000079929215000009/mho-20141231x10k.htm}{M/I Homes 10-k}
\end{singlespace}
\end{small}
\end{displayquote}


\begin{displayquote}
\begin{small}
\begin{singlespace}
``'In addition, liquidity requirements are dependent on the Company's credit ratings and general perception of its creditworthiness.''
\href{https://www.sec.gov/Archives/edgar/data/1013871/000101387114000003/a201310-k.htm}{NRG Energy 10-k}
\end{singlespace}
\end{small}
\end{displayquote}


\begin{displayquote}
\begin{small}
\begin{singlespace}
``PNMR, PNM, and TNMP cannot be sure that any of their current ratings will remain in effect for any given period of time or that a rating will not be put under review for a downgrade, lowered, or withdrawn entirely by a rating agency. Downgrades or changing requirements could result in increased borrowing costs due to higher interest rates in future financings, a smaller potential pool of investors, and decreased funding sources. Such conditions also could require the provision of additional support in the form of letters of credit and cash or other collateral to various counterparties.'' \href{https://www.sec.gov/Archives/edgar/data/22767/000110842614000016/pnm1231201310-k.htm}{PNM Resources 10-k}
\end{singlespace}
\end{small}
\end{displayquote}

\begin{displayquote}
\begin{small}
\begin{singlespace}

Under some over-the-counter derivative agreements on forms promulgated by the International Swaps and Derivatives Association, Inc. (“ISDA”), FGL has agreed to maintain certain financial strength ratings. A downgrade below these levels provides the counterparty under the agreement the right to terminate the open derivative contracts between the parties, at which time any amounts payable by FGL or the counterparty would be dependent on the market value of the underlying derivative contracts. FGL’s current rating allows multiple counterparties the right to terminate ISDA agreements, at which time the counterparty would unwind existing positions for fair market value. No ISDA agreements have been terminated, although the counterparties have reserved the right to terminate the ISDA agreements at any time. As of September 30, 2013, the amount at risk for ISDA agreements which could be terminated based upon FGL’s current ratings was \$221.8 million, which equals the fair value to FGL of the open over-the-counter call option positions. The fair value of the call options can never decrease below zero.'' \href{https://www.sec.gov/Archives/edgar/data/109177/000010917714000040/hrg-9302014x10xk.htm}{HRG Group 10-k}


\end{singlespace}
\end{small}
\end{displayquote}

\begin{displayquote}
\begin{small}
\begin{singlespace}


``Certain agreements with counterparties employ set-off, collateral support arrangements and other risk mitigating procedures designed to reduce the net exposure to credit risk in the event of counterparty default.'' \href{https://www.sec.gov/Archives/edgar/data/1467858/000146785814000043/gm201310k.htm}{General Motors 10-k}
\end{singlespace}
\end{small}
\end{displayquote}


\begin{displayquote}
\begin{small}
\begin{singlespace}
``Borrowings under the new credit facility bear interest at our option at the Eurodollar Rate (the LIBOR Rate) plus an applicable margin or the Base Rate (the highest of the Federal Funds Rate plus 0.50 \%, the 30-day Eurodollar Rate plus 1.0\%, or the administrative agent’s prime rate) plus an applicable margin. The applicable margins vary depending on our credit rating. Upon breach by us of certain covenants governing the new credit facility, amounts outstanding under the new credit facility, if any, may become due and payable immediately.'' \href{https://www.sec.gov/Archives/edgar/data/1179060/000117906014000006/xtex201310-k.htm}{Enlink Midstream Partners LP 10-k}
\end{singlespace}
\end{small}
\end{displayquote}

\begin{displayquote}
\begin{small}
\begin{singlespace}
``A downgrade of our credit rating might increase our cost of borrowing and could require us to post collateral with third parties, negatively impacting our available liquidity.'' \href{https://www.sec.gov/Archives/edgar/data/1338613/000133861314000012/rgp1231201310-k.htm}{Regency Energy Partners Lp 10-k}
\end{singlespace}
\end{small}
\end{displayquote}


\end{document}
